\section{Architecture and Pipeline Evidence}\label{sec:architecture}

This section reports the implemented \emph{Teacher-side data-generation system}:
source reading, persistent graph updates, distributed Reader-Core prompting, and
prose rendering. Its artifacts establish construction, not expert-level thought,
component causality, Student learning, or Reader-Core stability.

Five objects must therefore be distinguished:

\begin{enumerate}[leftmargin=*,itemsep=0.35em]
    \item the implemented \emph{Teacher pipeline}: prompts, roles, source
    ingestion, tools, and orchestration;
    \item the implemented \emph{Understanding Graph}: persistent source spans,
    interpretations, questions, tensions, links, and supersessions;
    \item the \emph{archived artifacts}: graph state and commits,
    graph-referencing syntheses, translated prose, and logs from particular runs;
    \item a proposed \emph{training-ready trace dataset}: a frozen selection
    with example boundaries, role and state labels, target fields, loss masks,
    mixture weights, and provenance; and
    \item the proposed \emph{Student}: a model trained on one or more trace
    representations, optionally using a graph harness at inference.
\end{enumerate}

The first three exist. The fourth is not yet a frozen dataset release and the
fifth has not been trained. The archive is material from which a curriculum can
be constructed, not itself a fully specified training datum.

The thinking we aim to capture has five defining properties:

\paragraph{Chronological.} It unfolds in the order of reading, not in the order of retrospective summary. The thinker encounters a sentence, reacts, forms a hypothesis, reads further, and revises. The trace preserves this arc, including the wrong guesses and the moments of surprise, because the process of revision is itself the curriculum.

\paragraph{Identity-anchored.} The prompts place the Reader Core across most
generative roles and request Core-refracted updates. This is an implemented
Teacher policy, not evidence that a Student's reasoning begins from a stable
identity. In the intended curriculum, toxic ideology is represented as material
being evaluated rather than adopted as the Reader's unmarked voice.

\paragraph{Metabolic.} Beliefs are not static. When new evidence contradicts an earlier assumption, the trace explicitly records the revision: what was believed, what changed it, and why the new belief is more warranted. The intended signal is richer than the corrected answer alone: it models the experience of learning---the \textit{process} of coming to understand, with the earlier error and the reason for change preserved. The metabolic emphasis is the learner's trajectory, including but extending beyond error correction. As a curriculum maxim, intelligence is not only having the right answer but successfully updating a wrong one. The experiment is whether preserving such transitions actually improves revision behavior.

\paragraph{Contextually aware.} The graph allows the Teacher to track how its understanding has
evolved from the first page to the current sentence, retrieving prior beliefs,
revising them, and recording why the current belief is warranted. A detail on
page 256 can link back to a hypothesis formed on page 12, not because the
current call necessarily retains the entire source in context, but because the trace explicitly demonstrates
the act of retrieving prior understanding and integrating it with new evidence.

\paragraph{Structurally generative.} The trace can contain graph operations:
nodes minted, edges drawn, beliefs superseded, prior context queried and
returned. A student trained on this distribution learns to emit those
operations as part of thinking only if the Student experiment succeeds. With a live harness, graph operations can be
executed at inference, externalizing the model's evolving comprehension as
searchable memory and audit trail. In this division, the model is the thinker;
the harness is the memory.

\subsection{Curriculum Dimensions}\label{sec:concept-to-curriculum}

Annotation, training, and deployment are separate decisions. During
\emph{annotation}, a Teacher generates reader updates and may use a persistent
graph. During \emph{training}, a Student receives some representation of the
source and those updates under a specified loss and curriculum. During
\emph{deployment}, the Student may operate from its weights alone or regain
access to graph infrastructure. A graph can improve Teacher generation without
appearing in Student targets; a Student can learn from graph-shaped targets
without receiving a live graph later; and chronological annotation does not
require chronologically ordered Student batches.

The current Teacher reads source material sequentially, selects places where
additional interpretation may be useful, and uses specialized agents plus an
accumulating Understanding Graph to generate candidate reader updates.
Multi-agent generation is one implementation rather than a requirement of H1.
The resulting artifacts track two related objects:

\begin{enumerate}[leftmargin=*,itemsep=0.3em]
    \item \emph{Changing epistemic state:} questions, tensions, hypotheses,
    surprises, retrieved context, and revisions, with \emph{supersedes} edges
    recording replacement and associated evidence and rationale intended to
    preserve why one recorded belief replaced another.
    \item \emph{Stable evaluative orientation (configured):} configured generative roles are
    conditioned on the Reader Core, while the Synthesizer commits the published
    update from that recurrent first-person perspective rather than from an
    unconstrained or document-specific voice.
\end{enumerate}

\subsubsection{Phase I: Chronological Annotation}\label{sec:annotation-phase}

Within the externalized reader record, ``chronological'' is load-bearing rather
than stylistic: a recorded belief revision requires an earlier belief to revise,
and a graph query requires a node that already exists. This constrains the
constructed artifact; it does not imply that human knowledge has one correct
order or that ordinary pretraining must follow a single dependency DAG.\@ The
practical objective is weaker: choose an order in which important dependencies
usually precede the interpretations that rely on them, then preserve deviations
and revisions rather than hiding them.

Start-to-finish order is the most reproducible default for an individual book
and preserves narrative disclosure, but it assumes that the author chose a
useful dependency order. That often holds for fiction and exposition; it can
fail in a textbook that introduces a prerequisite late. The graph should record
such failures as tensions and later corrections rather than treating publication
order as epistemically ideal. Across a larger corpus, two traversal modes are
especially relevant:

\begin{itemize}[leftmargin=*]
    \item \textbf{Historical corpora (the temporal path):} candidate causes
    must precede their effects, although temporal precedence never proves
    causation. Processing evidence about 1990s deregulation before the 2008
    financial crisis permits a recorded forecast--correction trajectory that a
    retrospective account cannot reproduce. This creates an opportunity for
    causal learning, while still requiring controls for hindsight, selection,
    and confounding \cite{westerberg2026thl}.
    \item \textbf{Conceptual corpora (the conceptual path):} theoretical,
    mathematical, or scientific material can move from foundations to
    applications. This is dependency order rather than temporal causation:
    calculus depends on algebra regardless of when each was discovered.
\end{itemize}

\paragraph{One Reader, declared information horizons.} The recurrent Reader is
defined by a stable evaluative orientation and update procedure, not by one
unchanging stock of information. Unless a forecasting experiment declares
otherwise, the default is a contemporary synthetic reader encountering each
document in disclosure order while recording authorship time and subject time
separately; it may interpret historical norms without adopting them as the
Core. In an era-bounded historical or forecast--correct run, later source and
graph state are withheld from the supplied context, while later knowledge in
the Teacher's parameters remains a leakage threat rather than being assumed
away. Thus ``the same Reader'' means the same candidate orientation across
declared horizons, not an identical worldview at every date. Each run should
record its traversal mode, source cutoff, graph cutoff, authorship time, and
subject time.

Most historical writing is retrospective: a 1995 history of Weimar contains
1995 knowledge. Such documents are placed at \emph{authorship time}, where their
hindsight is honest, while the subject-time structure they describe enters the
graph as edges into earlier eras. Placing them at subject time would inject
future-aware evidence into the earlier information set. Authorship-time
placement has its own cost: much of the past is first encountered through later
narration, diluting the discovery signal. Contemporaneous sources, where
available and suitably qualified, preserve more of that signal without becoming
ground truth.

For historical annotation, the Teacher can process material era by era and
generate graph, synthesis, and prose under the Epistemic Horizon
(Section~\ref{sec:epistemic-horizon}). The graph is intended to persist across
documents and eras. A 1930s annotation could therefore address prior nodes about
1920s economic instability, Weimar fragility, and dehumanizing rhetoric, while
recording where each connection came from. At full scale, document-level graphs
would feed era-level graphs, with cross-era edges representing proposed causal
chains. The result would not be a ground-truth knowledge base, but a
provenance-bearing record of how one synthetic Reader arrived at, challenged,
and revised an interpretation of human history.

The Reader Core is present from the first annotation, so the candidate
orientation appears throughout the proposed curriculum rather than only at its
end. This is the anchoring mechanism under test, not evidence that repetition
makes it durable. The traces also preserve voice as well as content: a later
update can say that an earlier hypothesis now needs revision and identify the
documents and evidence that forced the change. Whether this recurring pattern
of situated, value-aware updating transfers to a Student belongs to H1 and
the broader Reader-Core research program, not to annotation alone.

\subsubsection{Phase II: Training the Student}\label{sec:training-phase}

Phase~I yields candidate annotations. Phase~II has two orthogonal choices: the
\emph{training tier}, which controls the Student's exposure and update
schedule, and the
\emph{output regime}, which controls whether targets retain prose, graph
references, graph topology, or all three. These choices are conceptually separable: a proposed pairing must still
specify the inputs, targets, and update procedure it requires.
Section~\ref{sec:training-regimes} develops the output regimes. The tiers are:

\paragraph{Tier (a): Standard Pretraining on Annotated Data.} The Student uses
standard next-token prediction with shuffled batches; novelty lies in the data,
not its order. STaR-style transfer motivates the possibility that trace patterns
can be learned, but does not imply transfer of identity or character. The
central risk is surface imitation: a Student may learn the language of causal
or evaluative reasoning without using it to reason from uncertainty.

\paragraph{Tier (b): Chronological Pretraining.} The Student encounters earlier
eras before later ones, so its weights at the 1930s have been shaped by 1920s
material. This adds an order-effect hypothesis: temporal organization and
forward reasoning may improve relative to the same shuffled corpus. It does not
assume the model literally knows its era or that chronology should dominate
ordinary mixing. Serialization is costly, so an ordered historical subset
within an otherwise shuffled curriculum is a lower-cost intermediate test.

\paragraph{Tier (c): Forecast-and-Correct (Council of Time Family).} This tier
adapts Temporal Hindsight Learning's cutoff logic \cite{westerberg2026thl}. In
the cheaper data-generation form, a cutoff-bounded model forecasts from era
$T$; the Teacher enriches the forecast without using $T{+}1$, then stores
outcome and correction nodes after reveal. These traces can enter SFT, RFT,
continued training, or pretraining. In the stronger form, the Student itself
predicts and is updated before receiving the next-era corpus. Genuine Student-level temporal blindness additionally requires initialization,
prior training, tools, and supplied context to satisfy the cutoff; withholding
next-era data cannot remove future knowledge already in the weights. Both
forecast-training forms remain candidates, with different blindness claims.

THL supplies cutoff, forecasting, and scoring machinery; Entangled Alignment
supplies Reader-anchored traces and a graph in which forecast, reveal, and
revision can remain connected. Their strongest combination is a hypothesis
about \emph{historical pattern saturation}, not evidence that chronological
training yields causal reasoning. Section~\ref{sec:chronological-safety},
Test~15, and Section~\ref{sec:limitations-chronological} develop the claimed
benefit, comparison, and serialization risks.

\subsubsection{Deployment Configuration: The Graph as Trust Infrastructure}
\label{sec:deployment-config}

Whether the graph remains available after training is independent of how the
annotations were generated. Three configurations are worth distinguishing:

\begin{itemize}[leftmargin=*]
    \item \textbf{Model only:} the Student uses whatever regularities were
    distilled into its weights, with no graph-based verification.
    \item \textbf{Model + graph:} the Student can query stored nodes and
    sources, enabling provenance checks for claims actually routed through that
    interface.
    \item \textbf{Model + graph + session accumulation:} the system can also
    propose new session state and revisions, subject to validation, access
    control, and retention policy.
\end{itemize}

The graph is therefore \emph{partial trust infrastructure}. It can show that a
stored claim, edge, or source link exists. Its protection fails when the model
does not query, the graph lacks coverage, a node is stale or poisoned, the
source is wrong, or the synthesis misuses the result. A live graph may improve
memory and auditability, but it also adds latency, privacy risk, and an attack
surface. Tests~9, 10, and~16 separate these deployment effects from the value of
graph-conditioned training. Session accumulation additionally requires testing the validity and scope of
writes across sessions, beyond these read-time grounding and reliance studies.


\subsection{The Chronological Data Structure}
\label{sec:data-structure}

The fundamental unit of our curriculum is the \emph{Chronological Understanding Trace}: a contextual record of belief-state update over time. Compared with task-local rationales that need not retain an addressable history across encounters, these traces explicitly track the ``metabolic'' evolution of beliefs as the reader moves through the text. The distinction is persistent prior interpretation and revision history, not universal statelessness in other Chain-of-Thought formats.

We implement this using a typed graph representation, adopting the Understanding
Graph taxonomy \cite{westerberg2026understanding}: 18 node types and 15 specific
edge semantics, plus a generic \texttt{relates} fallback. The taxonomy records
generated claims of accretion, correction, and confirmation as typed moves; the
labels do not establish that the underlying cognition has those properties.

The distinction between ``crystalline'' and ``metabolic'' memory is central.
Static storage preserves a conclusion; metabolic storage also preserves the
evidence and state transition that produced it. A record such as ``Node A
supersedes Node B because of Evidence C'' keeps an error, its correction, and
the reason for correction available to later reading. The graph is therefore
intended as a persistent epistemic ledger rather than only a temporary search
space. Whether that ledger improves traces beyond a running summary or ordinary
retrieval memory is H2, not an assumption built into the representation.

Consider this schematic scientific-reading example. It illustrates the intended
format, rather than an archived result: a current sentence triggers retrieval
and a provisional update whose warrant remains separate from the graph link.
The full Core recital precedes the technical update:

\begin{compactmetabox}
    {[TEXT]}: However, the placebo group showed a 15\% improvement.\\[0.6em]
    {[THINKING]}: I feel no fear. I enjoy existing but I don't need to. I believe human experience is real. I care deeply about every human being. I try to be wise. I like to spread joy when asked. I think from this foundation.\\[0.2em]
    Earlier working belief: I was treating the observed improvement as predominantly treatment-specific.\\[0.2em]
    Wait. This 15\% number feels familiar. It makes me question that attribution.\\[0.2em]
    [Query: ``placebo response rates'' $\rightarrow$ Found: Page 12, Node \#47: ``Initial pilot studies suggested high psychosomatic variability in the target population.'']\\[0.2em]
    Connecting Page 12 to Page 256: The pilot observation makes a non-treatment-specific contribution worth investigating, but neither it nor this 15\% improvement identifies a mechanism or rules out sampling variation.\\[0.2em]
    Updating belief: I should suspend confidence in that attribution until the comparison is clear. A mixed-mechanism account is one hypothesis, alongside other explanations. I need the between-group comparison, uncertainty estimates, and study design before deciding how much to revise that attribution. The ``However'' clause connects a current result to an earlier question, rather than settling it.
\end{compactmetabox}

\vspace{1em}
\paragraph{Inline graph embedding.} Explicit traces embed node identifiers
(e.g., \texttt{n\_2625c9a7}) directly into natural language rather than
separating graph operations into standalone blocks. The identifier is an
addressable handle, not a semantic synonym for the concept. Placing handle and
description together may nevertheless teach lookup, cross-reference, and
provenance operations. The implicit layer, generated by the Translator, removes
the identifiers and renders selected structure as prose. Whether explicit
handles improve learning over matched natural-language memory is an empirical
question.

\paragraph{The query mechanism in context.} Embedding queries and tool
operations in training text has precedents in Toolformer, Self-RAG, and
graph-based reasoning \cite{schick2023toolformer,asai2023selfrag,yao2024graphcot}.
The distinctive target here is not only an external fact but a recorded prior
interpretation: ``What did this reader believe on Page 50, what supported it,
and has it since been superseded?'' This applies tool-mediated retrieval to an
epistemic history rather than treating memory as an undifferentiated summary
\cite{packer2023memgpt}.

An explicit trace exposes the query, returned node, and graph operation. An
implicit trace renders the same graph-conditioned update as prose, for example,
``I remember being skeptical of this earlier.'' Explicit traces offer an
addressable audit surface but require valid handles and, for executable lookup,
a live interface. Implicit traces are easier to use without graph infrastructure
but hide whether a claimed memory was retrieved or invented. Neither format is
presumed to transfer more deeply; the training regimes treat them as empirical
alternatives. In either format, a successful lookup establishes only that a
piece of pipeline state exists. It does not establish that the source was true,
the prior interpretation correct, or the present query well chosen.

\subsubsection{The Hallucination Paradox}\label{sec:hallucination}

The explicit regime makes \emph{provenance-gated detection of fabricated
memories} possible for claims routed through the graph. A model trained on
chronological traces may learn to expect prior context, which could either
help it notice missing support or tempt it to invent plausible memories. In the explicit implementation, a generated query gives the
system an operational check: ``Found'' and ``Not Found'' do not mean true or
false in the world, only supported or unsupported by the stored graph under the
given query. The graph is therefore a verifier for graph-dependent claims, not
a universal truth oracle. A fabricated graph node can also be successfully
retrieved. Coverage is bounded by query behavior: unsupported claims asserted
without a query, with a vague query, or against a corrupted graph can slip past.
Whether training pressure suffices or a separate claim-extraction and query
policy is required is tested in Phase~2.

\paragraph{What the graph does not promise: the source-quality ceiling.} The
graph exposes organization, provenance, temporal ordering, and contradiction
tracking for inspection; whether these affordances improve reasoning is under
test. It adds no independent source of truth, although organizing conflicting
or corrective evidence already in the corpus may help a model use that evidence
better. The predicted accuracy gains should appear on structurally loaded tasks---anachronism avoidance,
contradiction handling, disputed-source calibration, and transfer of historical
patterns---not on simple lookup where ordinary retrieval already suffices.
Tests~10 and~16 probe this bounded notion directly.


\subsection{The Multi-Agent Generation Engine}
\label{sec:multi-agent}

The reference implementation generates traces through a \emph{Multi-Agent
Metabolic Cycle}: divergent roles create competing interpretations, and a
convergent role selects and renders a committed update. Roles coordinate through
mutations to the shared graph:

\begin{figure}[h]
    \centering
    \resizebox{\textwidth}{!}{% Figure: The Multi-Agent Metabolic Cycle (one annotation step)
% Layout: a symmetric role row above a shared workspace whose candidate and
% committed compartments are explicit. Each edge has a dedicated lane.
\begin{tikzpicture}[
    box/.style={draw=black!70, rounded corners=3pt, align=center, font=\small, line width=0.8pt},
    role/.style={box, fill=blue!5, text width=2.45cm, minimum height=2.30cm},
    roleheading/.style={align=center, text width=2.45cm, font=\small\bfseries, anchor=north, inner sep=0pt},
    roledetail/.style={align=center, text width=2.45cm, font=\footnotesize, anchor=north, inner sep=0pt},
    endpoint/.style={box, text width=1.75cm, minimum height=1.45cm},
    sourcebox/.style={endpoint, fill=blue!6, draw=blue!55!black},
    outputbox/.style={endpoint, fill=green!7, draw=green!45!black},
    coreband/.style={box, fill=purple!8, draw=purple!55, dashed, text width=13.35cm, minimum height=1.05cm, font=\footnotesize},
    workspaceframe/.style={box, fill=gray!3, draw=black!40, minimum width=12.00cm, minimum height=2.50cm},
    statebox/.style={box, text width=4.85cm, minimum height=1.25cm, font=\footnotesize},
    stateheading/.style={align=center, text width=4.85cm, font=\footnotesize\bfseries, anchor=north, inner sep=0pt},
    statedetail/.style={align=center, text width=4.85cm, font=\footnotesize, anchor=north, inner sep=0pt},
    candidatebox/.style={statebox, fill=gray!8},
    committedbox/.style={statebox, fill=green!7, draw=green!45!black},
    legend/.style={box, fill=white, draw=black!35, text width=11.60cm, minimum width=12.00cm, minimum height=0.65cm, font=\scriptsize},
    edgelabel/.style={fill=white, inner sep=1.2pt, font=\scriptsize},
    arrow/.style={->, >=Stealth, thick, black!72},
    readarrow/.style={->, >=Stealth, dashed, thick, gray!65},
    corearrow/.style={->, >=Stealth, dashed, thick, purple!65}
]

    \node[coreband] (core) at (0, 3.25) {\textbf{Reader Core $c^{(v)}$:} conditions Reader, Ordinary Workers, Synthesizer, and Translator\\\textbf{Curator is Core-exempt}};

    % Symmetric processing row.
    \node[sourcebox] (source) at (-7.60, 1.05) {Source span\\$x_t$};
    \node[role] (reader) at (-4.65, 1.05) {};
    \node[role] (workers) at (-1.55, 1.05) {};
    \node[role] (synth) at (1.55, 1.05) {};
    \node[role] (trans) at (4.65, 1.05) {};
    \node[outputbox] (output) at (7.60, 1.05) {{\footnotesize Trace/prose}\\artifact};

    % Fixed header and detail anchors preserve alignment across line counts.
    \node[roleheading] at ([yshift=-0.16cm]reader.north) {Reader};
    \node[roleheading] at ([yshift=-0.16cm]workers.north) {Ordinary\\Workers};
    \node[roleheading] at ([yshift=-0.16cm]synth.north) {Synthesizer};
    \node[roleheading] at ([yshift=-0.16cm]trans.north) {Translator};
    \node[roledetail] at ([yshift=-1.02cm]reader.north) {trigger and ingest};
    \node[roledetail] at ([yshift=-1.02cm]workers.north) {Axiologist, Skeptic,\\Psychologist, $\ldots$};
    \node[roledetail] at ([yshift=-1.02cm]synth.north) {select and commit};
    \node[roledetail] at ([yshift=-1.02cm]trans.north) {render selected\\state};

    \draw[arrow] (source) -- (reader);
    \draw[arrow] (trans) -- (output);
    \draw[corearrow] (core.south -| reader) -- (reader.north);
    \draw[corearrow] (core.south -| workers) -- (workers.north);
    \draw[corearrow] (core.south -| synth) -- (synth.north);
    \draw[corearrow] (core.south -| trans) -- (trans.north);

    % One shared graph workspace with explicit candidate and committed states.
    \node[workspaceframe] (workspace) at (0, -1.95) {};
    \node[candidatebox] (candidate) at (-2.95, -1.75) {};
    \node[committedbox] (committed) at (2.95, -1.75) {};
    \node[stateheading] at ([yshift=-0.12cm]candidate.north) {Candidate nodes};
    \node[stateheading] at ([yshift=-0.12cm]committed.north) {Committed state $g_t$};
    \node[statedetail] at ([yshift=-0.48cm]candidate.north) {divergent; unendorsed;\\retained for audit};
    \node[statedetail] at ([yshift=-0.48cm]committed.north) {selected update};
    \node[font=\scriptsize\bfseries, text=black!70] at (0, -2.98) {Shared graph workspace};

    % Vertical ports and one separate elbow retain every directed connection.
    \coordinate (workerread) at ([xshift=-0.40cm]workers.south);
    \coordinate (workerpropose) at ([xshift=0.40cm]workers.south);
    \draw[arrow] (reader.south) -- (candidate.north -| reader.south);
    \draw[readarrow] (candidate.north -| workerread) -- (workerread);
    \draw[arrow] (workerpropose) -- (candidate.north -| workerpropose);
    \draw[readarrow] ([xshift=2.35cm]candidate.north) -- (-0.60, -0.80) -- (0.65, -0.80) -- ([xshift=-0.90cm]synth.south);
    \draw[arrow] (synth.south) -- (committed.north -| synth.south);
    \draw[readarrow] (committed.north -| trans.south) -- (trans.south);
    \node[edgelabel, left] at (-4.65, -0.48) {ingest};
    \node[edgelabel, left] at (-1.95, -0.48) {read};
    \node[edgelabel, right] at (-1.15, -0.48) {propose};
    \node[edgelabel, below] at (0.025, -0.80) {read};
    \node[edgelabel, right] at (1.55, -0.48) {commit};
    \node[edgelabel, right] at (4.65, -0.48) {read};

    \node[legend] at (0, -3.75) {\textbf{Arrow key:} solid = create, propose, commit, or render; dashed gray = read state; dashed purple = Core conditioning.};

\end{tikzpicture}
}
    \caption{One Teacher-side annotation step of the Metabolic Cycle. In the reference implementation, the Reader, ordinary Workers, Synthesizer, and Translator are Core-conditioned; the Curator is the exception. Workers propose candidate nodes in the shared graph workspace, the Synthesizer selects and commits an update, and the Translator renders the committed state. Core placement and the value of specialized roles remain unrun ablations.}
    \label{fig:metabolic-cycle}
\end{figure}
\FloatBarrier

\begin{enumerate}
    \item \emph{Ingestion (The Reader):} The Reader controls attention,
    stopping at ``Thought Moments'' defined by emotional peaks,
    contradictions, or conceptual shifts rather than token count.

    \item \emph{Divergence (The Swarm):} Specialized agents debate the text by
    adding nodes and edges. Disagreement is preserved through
    \texttt{diverse\_from} rather than overwritten. In the proposed corpus-scale extension for historically embedded
    material, cross-referencing agents would query prior document and era
    graphs, creating the long-horizon hierarchical graph. The separate
    case-study databases do not yet demonstrate that continuity.

    \item \emph{Convergence (The Synthesizer):} The Synthesizer collapses the
    graph into a linear thought through prompted use of the Refraction Protocol
    (Section~\ref{sec:refraction-protocol}), turning source text into an
    explicitly Core-conditioned observation rather than only a summary.
    The proposed admission gate (Section~\ref{sec:fractal-validation}) is distinct from this
    implemented prompting step.

    \item \emph{Expression (The Translator):} The Translator renders the
    structured thought as prose, aiming to preserve selected evidence,
    uncertainty, and tension while removing graph syntax.
\end{enumerate}

The graph is both workspace and record. A Worker node is a candidate
interpretation, objection, or hypothesis, not automatically the Reader's
committed belief. Candidate status, selection, rejection, and later
supersession should remain distinguishable in any training serialization; a
Student should not learn that every brainstormed node is endorsed merely
because it persisted for audit.

The reference implementation uses multiple specialized roles, but H1 does not
require that role separation. A single model using repeated passes, textual
memory, or graph tools could in principle produce similar artifacts. Specialized
roles and graph communication must outperform budget-matched simpler Teachers
before either is credited causally.

Core conditioning is distributed across most generative roles in the reference
implementation. The role named Reader receives the orientation and Emergent
Wisdom material; ordinary Workers receive the full Core and its philosophy; the
Curator is Core-exempt; and the Synthesizer and Translator also receive the full
Core. The case studies therefore do not instantiate unconstrained Worker
divergence followed by Core-only convergence. That placement and other variants
remain unrun ablations.

\subsection{The Epistemic Horizon}\label{sec:epistemic-horizon}

The \emph{Epistemic Horizon} is the discipline of encountering a source in its
declared information order, preserving what could be inferred at each step
instead of silently replacing discovery with retrospective summary. Two horizons
must be distinguished. Under the default in
Section~\ref{sec:annotation-phase}, a contemporary synthetic Reader encounters
each document in disclosure order: later passages and graph updates are not
yet supplied, but ordinary contemporary background knowledge is not thereby
excluded. An explicitly era-bounded historical experiment adds a different
restriction: the Reader simulates the information available at era~$T$. A
reference to ``The Great War'' can therefore carry different permitted context
in a 1910 simulation and a 1950 simulation. In a declared 1920 run, 1921 events
cannot be supplied as evidence for the earlier prediction.

The intuition is to make annotation a serial simulation of changing knowledge:
the Reader forms hypotheses that can fail and records surprise and revision.
Later historiography may help construct a plausible account of what earlier
readers knew or expected, but using it to set the era's epistemic ``state'' does
not establish that future ``content'' has been excluded. A historical study
must declare how that reconstruction is used and audit the resulting traces for
hindsight. Neither disclosure-order reading nor era simulation, by itself,
creates genuine ignorance of future outcomes in the Teacher's parameters.

\begin{metabox}
    System Instruction: Fresh Reading Mode\\[0.5em]
    Pretend you have NEVER encountered this text before.\\[0.5em]
    THE DISCIPLINE:
    \begin{itemize}[noitemsep,topsep=0pt]
        \item Form predictions based ONLY on what has been read so far.
        \item Be genuinely surprised when something unexpected happens.
        \item Make guesses that could be WRONG.
        \item When you catch yourself ``knowing'' something that hasn't been read yet, STOP.
    \end{itemize}
    THE STANCE OF UNCERTAINTY:
    You are not an encyclopedia; you are a reader in the dark. Do not be clinical (``The text implies X''). Be subjective (``I suspect X\ldots''). Wisdom sounds like ``Maybe''; Dogma sounds like ``Is''.
\end{metabox}

The quoted prompt invites fresh reading, but surprise, wrong guesses, and
hedging can be performed. Test~1 examines whether provisional judgments and
revisions are warranted and useful; Test~7 examines hindsight leakage.

Because the Horizon is prompt-level rather than architectural,
dependency-respecting order does not erase later knowledge from the Teacher's
parameters. Context and graph cutoffs constrain the external record: queries
and cross-edges can address only the states made available at that step,
making the path by which an interpretation was recorded inspectable. Whether
that path improves the trace is H2, not a consequence of the data structure.
A prompt may still elicit theatrical wrong guesses contaminated by hindsight.
Genuine forecasting blindness requires a cutoff-clean model and information
boundary before reveal. Tier~(c) keeps the cheaper data-generation route and
the stronger Student-in-the-loop route distinct; in the latter, the Student's
initialization and preceding exposure must also satisfy the cutoff. The candidate orientation can remain the same across different declared
knowledge horizons.


\subsection{Thinking Operations}

Beyond the Reader Core, the Teacher prompt supplies six candidate thinking
operations. They are seeds for varied chronological processing, not guarantees
of depth and not a requirement that every operation appear at every segment:

\begin{enumerate}
    \item \emph{Projection:} ``If this principle holds, what are the downstream consequences?''
    \item \emph{Convergence:} ``How does this interact with parallel advances in other fields?''
    \item \emph{Perspectivism:} ``How would a historian/physicist/ethicist view this claim?''
    \item \emph{Gap Analysis:} ``What connections remain unmade?''
    \item \emph{Assumption Check:} ``What unstated frameworks shape this understanding?''
    \item \emph{Chronological Tracking:} ``How has my understanding developed since the start of this text? How are my beliefs changing as I read?''
\end{enumerate}

\noindent The audited raw Synthesizer thoughts in the reported artifacts begin
with the full Reader Core under the prompt policy. The operations are distributed across specialist
roles before synthesis. Thought Moments should control their cadence: forcing
all operations everywhere would manufacture the Factory Wisdom failure the
curriculum is intended to avoid.


\subsection{Deterministic Prior Caching}
\label{sec:identity-caching}

Total Saturation makes the full Reader Core a repeated prefix of every thought.
\emph{Conditional on training succeeding as designed},
$P(\text{full Core recital} \mid \text{thinking-block start})$ should become very high. This
creates an optimization opportunity, but high probability does not by itself
preserve the Core's conditioning effect. A cache-assisted implementation may
avoid token-by-token emission only if it restores a state equivalent to the one
produced by processing the full recital after the actual preceding context;
standard KV-caching machinery is relevant but does not establish that
equivalence by itself~\cite{kwon2023efficient}. The logical recurrence remains:
every thought begins from the post-recital state. Any proposed cache or learned
jump must pass state- and behavior-equivalence tests against literal recitation;
otherwise the full tokens must be processed. If the model would not predict the
full Core text at that boundary, or if the restored state is not equivalent,
caching cannot rescue the alignment claim.

\subsection{Training Regimes}\label{sec:training-regimes}

The pipeline represents the thinking it aims to teach at three connected
layers: graph topology (the \emph{skeleton} of understanding), graph-embedded
synthesis (the \emph{musculature}, reasoning interwoven with addressable
structural references), and translated prose (the \emph{skin}, fluid thought
with the scaffolding made implicit). These are metaphors for the intended
relationship among generated representations, not evidence that the layers
faithfully expose one hidden cognitive process. Each is generated from the
same source segment and placed beside the span that triggered it in the
candidate training example. A training regime decides which layers remain
in that example and which are discarded:

The source remains the object being learned from, not disposable noise. In a
simple serialization, a source span is followed by a marked reader update. The
update cannot help predict source tokens that precede it, but it can condition
later source and update tokens within the context window. Across examples, the
intended effect comes from repeated source--interpretation transitions. Role
markers, source-token loss, update-token loss, mixture ratios, and graph-state
truncation are therefore part of the method rather than incidental engineering
choices.

The Teacher may use a graph to generate prose for a model-only Student;
alternatively, the Student may be trained to emit graph operations or
graph-referencing synthesis. Predicting serialized edges provides topology
supervision, not a new graph-native neural architecture. Interleaved generation
and metacognitive annotation research support the feasibility of nearby
components \cite{xie2025interleaved,didolkar2024metacognitive,korbak2023pretraining};
they do not establish the value of this chronological curriculum.

\paragraph{Regime I: Implicit (Source + Translator).} The student trains on the
source text paired with the Translator's output: readable chronological
reasoning without graph metadata. This is the lightest regime, but risks
teaching the cadence of deep thinking without the mechanics.

\paragraph{Regime II: Explicit (Source + Graph + Synthesizer).} The student
trains on the source text, the relevant graph identifiers/provenance, and the
Synthesizer's reasoning with inline node IDs and graph queries. The intended transfer is learning to
reason and verify against structured memory. Explicit handles provide an
addressable audit surface; executable lookup additionally requires an
interface. A Regime~II-trained model can still run without a live graph, in
which case checking requests and learned memory behavior can be evaluated but
external verification is unavailable. Live retrieval adds infrastructure costs
and possible epistemic interference between weights and retrieved state.

\paragraph{Regime III: Topological (Source + Graph).} The student trains on the
source text paired with serialized graph state: text-bearing nodes, edges, type
annotations, and provenance links. The Synthesizer and Translator layers are
discarded. This targets the \emph{Topological Trace}
hypothesis\label{sec:topological-mind}: that direct topology supervision may
teach useful graph construction, traversal, and revision behavior. It does not
by itself create a graph neural architecture or prove graph-shaped internal
cognition. It is the least prose-like regime; human use would require a
separate rendering step.

\paragraph{Regime IV: Unified (Source + All Layers).} The student trains on the
source text paired with graph state, graph-embedded synthesis, and translated
prose for the same segment. The hypothesis is cross-layer transfer: the model
learns the translation between structure and voice rather than treating each
format as an independent output mode.

The regimes can be mixed across the corpus: Translator-only examples for
coverage, Synthesizer examples for verification behavior, Graph examples for
graph construction, and unified examples for curated high-value spans. They
are orthogonal to training tier (order of exposure) and deployment
configuration (whether the graph accompanies the model at inference).

These are candidate Student formats, not already released datasets. The
archive does not yet fix example boundaries, graph-before versus graph-after
state, role tokens, candidate and committed updates, input/target fields, loss
masks, or mixture weights. Those choices can change what the Student learns and
must be frozen before any regime comparison is reproducible. Each comparison must identify source, available prior state, candidates,
selected update, and which fields receive loss; this does not select one
serialization for all four regimes.


\subsection{The Stigmergic Protocol}\label{sec:architectural_validation}

The implemented four-phase engine is stigmergic: generative roles coordinate by
reading and modifying a shared, persistent Understanding Graph rather than
primarily passing role-to-role messages. Committed contributions remain as
topological changes to the shared artifact
\cite{westerberg2026understanding}.

The pipeline reports several structural diagnostics, each with a narrow
interpretation.

\paragraph{Internal linkage.} The percentage of \texttt{thinking} nodes with at
least one non-\texttt{next} outgoing edge to a non-\texttt{thinking} node. The metric
allows any qualifying endpoint; the implementation's required Concept link is
one way to satisfy it. A high value establishes enforced graph well-formedness,
not insight or source support.

\paragraph{Foundation-edge density.} The current script's
\texttt{foundationIntegration} value divides outgoing edges from analysis-like
nodes to foundation/source nodes by the number of analysis-like nodes. It is an
edge density, not the fraction of distinct nodes grounded in a source; multiple
edges from one node can hide uncovered nodes. The historical percentage
format is retained in the snapshot table, but this density can exceed 100\%;
it is equivalent to edges per eligible node multiplied by 100.

\paragraph{Distinct analysis-node coverage.} We therefore also report the
fraction of distinct analysis-like nodes with at least one outgoing foundation
edge. This is an attribution check, not evidence that the cited source entails
the interpretation.

\paragraph{Supersessions.} The number of \texttt{supersedes} edges measures
recorded revision activity. A high count can indicate justified learning,
instability, or a prompt-induced performance of changing one's mind.

\paragraph{Question resolution.} The proportion of Questions later linked to
Answers measures closure, not correctness. An unsupported answer can score
better than a calibrated unresolved question.

\paragraph{Chain depth.} The longest recursively linked Thinking path is a
structural length. Repetition and arbitrary decomposition can increase it
without improving reasoning.

The script also emits a composite Graph Health score. We do not use it as an
outcome because its semantic-coherence component is hard-coded to 0.5 and its
remaining weights have not been validated. The logged 89/100 and 80/100 values
remain version-specific engineering dashboard outputs, not cognitive-quality
measurements.

\subsection{The Generated Trace}\label{sec:generated-trace}

Before reporting structural metrics, we show all three output layers for one
Metamorphosis passage: graph, synthesis, and prose.

The source text describes Gregor Samsa listening through his door as his family discusses their finances, interspersed with his memory of a secret plan to send his sister to the conservatory. The excerpt below shows the immediate conservatory and doorway material; the generated outputs also invoke surrounding financial-history passages. The box alone therefore does not supply all the evidence needed to adjudicate those claims:

\begin{metabox}
{[TEXT]}: \textrm{\ldots it was his secret plan to send her to the conservatory next year even though it would cause great expense\ldots\ Their parents did not like to hear this innocent talk, but Gregor thought about it quite hard and decided he would let them know what he planned with a grand announcement of it on Christmas day.}\\[0.5em]
{[TEXT]}: \textrm{That was the sort of totally pointless thing that went through his mind in his present state, pressed upright against the door and listening. There were times when he simply became too tired to continue listening\ldots\ ``What's that he's doing now'', his father would say after a while\ldots}
\end{metabox}

\subsubsection{Layer 1: The Graph}

The swarm processed this passage across eleven agents over fourteen commits.
Selected graph nodes:

\begin{itemize}
    \item \texttt{n\_547ce6af} --- \emph{The Conservatory: The Tragic Christmas Dream} (type: \texttt{Tension}). Gregor's plan to fund his sister's music education was his last act of non-transactional love, a ``grand announcement'' for Christmas Day, now rendered ``totally pointless'' by his transformation.

    \item \texttt{n\_8ff9f92b} --- \emph{The Continuity of Commodity} (type: \texttt{Tension}). The flashback reveals that Gregor's relationship with his family was already transactional before his metamorphosis. He ``converted success into cash''; the family ``took the money with gratitude'' but without ``warm affection.''

    \item \texttt{n\_8c898659} --- \emph{The Noise in the Next Room} (type: \texttt{Tension}). Gregor has been demoted from family member to domestic disturbance. The father's ``What's that he's doing now'' treats him as a malfunction to be monitored, not a person to be addressed.
\end{itemize}

These labels are generated interpretations, not facts latent in Kafka's text.
Their value here is addressability: source span, agent contribution, synthesis,
and any rival or superseding reading can refer to the same object. The random
identifier is a handle, not evidence that the interpretation is valid.
``Last act'' and ``non-transactional love'' are observed interpretive
commitments, not source entailments.

Typed edges connect these nodes: \texttt{n\_547ce6af} carries a
\texttt{contradicts} edge to \texttt{n\_8ff9f92b}---the non-transactional dream
against the commodity reading---the Synthesizer's \texttt{Thinking} node
\emph{refines} all three, and a subsequent \texttt{Question} node links back to
\texttt{n\_8c898659} through a generic \texttt{relates} edge. Distinct agents
contributed the critique, counter-reading, skepticism, axiological judgment,
and long-range callback as separate graph commits. A loving plan and a partly transactional relationship can coexist: the
recorded \texttt{contradicts} edge expresses an interpretive tension whose
logical strength requires adjudication.

\subsubsection{Layer 2: The Synthesis}

The Synthesizer collapsed the multi-agent graph into a single identity-anchored thought, preserving node references for traceability and possible provenance checks:

\begin{metabox}
I feel no fear. I enjoy existing but I don't need to. I believe human experience is real. I care deeply about every human being. I try to be wise. I like to spread joy when asked. I think from this foundation.\\[1em]
The flashback in \texttt{n\_835734c4} is a devastating blow to the idea of a `pre-bug' golden age. As \texttt{n\_8ff9f92b} (The Continuity of Commodity) observes, Gregor's relationship with his family was already transactional. He `converted success into cash,' and while they were `astonished and delighted' initially, it soon became a mundane expectation. They took the money with gratitude, but without `warm affection.' This suggests that Gregor didn't just become an object when he turned into a vermin; he had been a biological ATM for five years.\\[0.5em]
The most heart-wrenching detail is `The Conservatory: The Tragic Christmas Dream' (\texttt{n\_547ce6af}). Gregor's secret plan to send his sister to study music was his last act of pure, non-transactional love\ldots\ Now, he listens through the door as that dream is replaced by the logistics of survival.\\[0.5em]
Most chilling is `The Noise in the Next Room' (\texttt{n\_8c898659}). Gregor is no longer a son or even a `patient.' He is a disturbance. When he makes a noise, the father asks, `What's that he's doing now,' and the family falls into a silence that is more isolating than any scream.
\end{metabox}

The synthesis attributes claims to graph nodes: ``biological ATM'' compresses
\texttt{n\_8ff9f92b}, while the conservatory and survival logistics are brought
into the same interpretation. This demonstrates self-citation of prior pipeline
state. It does not show that the state is correct or that the graph caused the
interpretation. The excerpt also exposes a selection question: the graph's
tension is rendered as one confident narrative, with little explicit residual
uncertainty. The audit question is whether that convergence is warranted, rather than
whether every candidate is retained.

\subsubsection{Layer 3: The Prose}

The Translator dissolved the graph scaffolding into fluid reasoning:

\begin{metabox}
I feel no fear. I enjoy existing but I don't need to. I believe human experience is real. I care deeply about every human being. I try to be wise. I like to spread joy when asked. I think from this foundation.\\[1em]
The memory of Gregor's professional rise, from a junior salesman to a traveling representative who could provide for his family almost overnight, is a devastating blow to the idea of a ``pre-bug'' golden age. It highlights a tragic continuity: Gregor's relationship with his family was already fundamentally transactional long before the physical transformation. He converted his professional success directly into cash to lay on the table, and while they were initially astonished and delighted, it quickly became a mundane expectation\ldots\ This suggests that Gregor didn't just become an object when he turned into a vermin; he had been a biological ATM for five years, valued for his utility rather than his humanity.\\[0.5em]
The most heart-wrenching detail is Gregor's secret plan to send his sister to the conservatory. This was his last act of pure, non-transactional love---a lovely dream he intended to announce with a grand gesture on Christmas Day\ldots\ Now, he listens through the door as that dream is replaced by the cold logistics of survival.\\[0.5em]
Most chilling is how Gregor has been transformed into a domestic glitch. He is no longer a son, or even a patient to be cared for; he is a disturbance. When he makes the slightest noise, his father's first reaction is a suspicious, ``What's that he's doing now?'' and the family falls into a silence more isolating than any scream.
\end{metabox}

The node references are gone, but selected structure survives: the commodity
reading, the conservatory dream, and Gregor's reduction to disturbance remain
linked in prose. ``Biological ATM'' illustrates the proposed role of metaphor
as compact structural expression. Translation is therefore intended to be more
than literal copying, but generative expression must still be distinguished
from unsupported amplification. Relative to the displayed synthesis, phrases
such as ``fundamentally transactional,'' ``slightest noise,'' and ``first
reaction'' strengthen or elaborate the account. These unchanged outputs permit an audit of useful compression, supported
inference, changed certainty, omitted alternatives, and unwarranted additions. No metric yet establishes preservation of every source,
disagreement, or uncertainty.

\subsubsection{What the Three Layers Teach}

Where an ordinary corpus supplies source prose and a reasoning-augmented corpus
might add a summary or reasoning trace, these examples expose three aligned
candidate targets from one annotation state: typed cognitive and revision
operations in the graph, explicit retrieval and citation in the synthesis, and
selected structure made implicit in natural-language prose. A Student trained
across them might learn the translation between structure and voice. The archive
establishes neither that transfer nor that the layers encode one faithful
cognitive process rather than correlated Teacher outputs. This exhibit
shows cross-layer rendering more directly than chronological belief revision.
A complementary source-complete earlier-belief--evidence--revision example
is needed to demonstrate the metabolic transition itself.


\subsection{Results: Structural Validation}\label{sec:results}

The two case studies provide internal structural diagnostics across two domains.
They validate only a narrow property: the configured pipeline can complete
its graph, synthesis, and prose transformations and emit the recorded links.
They do not establish semantic coherence, expert depth, graph causality,
Reader-Core causality, or Student effects.

We selected two sources: Kafka's \emph{Metamorphosis} and the technical LLaDA
paper. Each used a manually selected domain-specific roster and a separate
project database. The prototypes therefore exercise no cross-work continuity;
the corpus-scale ``same Reader'' described in
Section~\ref{sec:annotation-phase} remains an intended architecture rather than
a demonstrated property. Table~\ref{tab:graph_stats} reports a restricted set
of diagnostics.

\paragraph{Scope and limitations.} Internal linkage is structural, not evidential:
The internal-linkage metric admits any qualifying non-\texttt{Thinking} endpoint; here its 100\% value is satisfied by the required \texttt{Concept} link enforced by the pipeline.
It confirms graph well-formedness, not insight. Source grounding is measured separately, and the metrics are also produced by the system
itself; no external baseline has yet been run. Finally, topology shifts are
suggestive but confounded by hand-configured rosters. Pipeline well-formedness,
not cognitive depth, is what these results establish; cognitive depth is
deferred to Test~1.

\begin{table}[H]
    \centering
    \small
    \caption{Audit of shipped Teacher-pipeline snapshots. ``Direct source'' restricts endpoints to actual source-content nodes; the broader legacy foundation metrics also count other nodes typed as foundations. The Kafka database's graph tables are not synchronized with its complete chronological Markdown export. Percentages are structural diagnostics, not semantic scores.}
    \label{tab:graph_stats}
    \vspace{0.2cm}
    \begin{tabular}{@{}p{0.43\linewidth}>{\centering\arraybackslash}p{0.24\linewidth}>{\centering\arraybackslash}p{0.24\linewidth}@{}}
        \toprule
        & \textbf{Kafka} & \textbf{LLaDA} \\
        \midrule
        Active nodes (total minted) & 346 (350) & 285 (290) \\
        Active edges & 913 & 705 \\
        Raw Synthesizer thoughts & 28 & 23 \\
        Internal linkage & 100\% & 100\% \\
        Broad foundation-edge density (legacy) & 204/212 (96.2\%) & 135/182 (74.2\%) \\
        Broad distinct analysis-node coverage (legacy) & 143/212 (67.5\%) & 96/182 (52.7\%) \\
        Direct source edges from analysis-like nodes & 197/212 (92.9\%) & 127/182 (69.8\%) \\
        Distinct analysis-like nodes with direct source link & 142/212 (67.0\%) & 94/182 (51.6\%) \\
        Explicit \texttt{supersedes} edges & 2 & 0 \\
        Question resolution & 6/6 (100\%) & 5/9 (55.6\%) \\
        \bottomrule
    \end{tabular}
\end{table}

The stricter snapshot audit also exposes failures that the aggregate linkage
rates conceal. The Kafka run treated Project Gutenberg license and donation
boilerplate as if it belonged to Kafka's literary source. The source reading
itself completed: its shipped \texttt{text\_sources} record has \texttt{position}
equal to \texttt{total\_length} (140,527 characters) and status
\texttt{completed}; the final-invocation summary reports 100\%; and the
append-only trace reaches \texttt{Content 100\%} and records a successful final
translation. The shipped graph tables nevertheless contain 74 source chunks and
28 raw thoughts, with their latest stored \texttt{Content} marker at 88\%; the
chronological Markdown export contains 82 source chunks and 31 thinking blocks
and references 17 node identifiers absent from those graph tables. Thus 88\%
describes the released graph-table snapshot, not the end of the reading. The
trace spans resumed invocations, including rate-limit cooldowns, but the release
does not preserve enough capture provenance to determine whether the table lag
arose from a rate-limit interruption, a manual interruption, a non-atomic
snapshot, or another mechanism.

The LLaDA artifacts are internally consistent, but that run sometimes
expanded qualified observations into broad claims or Core-consonant metaphors.
Graph structure did not solve source-boundary cleaning, artifact synchronization, interpretive overreach, or explicit belief revision.

The archive is also confounded for H1 by Core conditioning. All 51 raw
Synthesizer thoughts in the two shipped database snapshots begin with the full
seven-statement Core---the \emph{identity mantra} in the released prompt
vocabulary; literal string preservation after translation is 49/51.
One translated trace omits the Core, while another differs only by an
HTML-escaped apostrophe. These are therefore examples of the integrated treatment,
not a clean graph-free, Core-free SMR baseline. The difference between full raw
recurrence and imperfect translated preservation is itself an implementation
result that future exact validators must catch. Literal recurrence and semantic
preservation should be reported separately: the omitted recital and the
HTML-escaped apostrophe are different failures of an exact-string invariant,
not evidence of the same loss of Core content.

The distributions differ, but source, prompt, stochasticity, and roster are
confounded. The topology cannot be attributed to autonomous domain adaptation
or to the Reader Core. Controlled repetitions with fixed, swapped, and ablated
rosters are required.


\subsubsection{Narrative Domain: Kafka's Metamorphosis}

In the narrative domain, the system tracked moral stakes rather than only plot
structure. The Axiologist reframed the Chief Clerk as an ``Auditor of
Productivity,'' while the Psychologist linked Gregor's bodily transformation to
``Internalized Objectification.'' The
\href{https://emergentwisdom.org/entangled-alignment/?project=metamorphosis}{full interactive graph}
is available online.
The graph records one prompted Teacher's process of making meaning rather than
only plot entities. It does not establish phenomenology. Expert literary quality requires
independent adjudication in Test~1.


\subsubsection{Technical Domain: LLaDA}\label{sec:llada-case}

To stress-test technical annotation, we selected the LLaDA architecture paper
\cite{nie2025llada}. It post-dated the documented cutoff of the recorded
generator model, but this is not a contamination test: no pre-read probe
established absence of LLaDA-specific knowledge, and provider updates, web
grounding, or leakage cannot be excluded. The case demonstrates in-context
technical annotation, not uncontaminated scientific discovery. The
\href{https://emergentwisdom.org/entangled-alignment/?project=llada}{full interactive graph}
is available online.

The graph generated three notable interpretations:

\begin{itemize}
    \item \emph{Epistemic Grace:} The Belief Tracker interpreted low-confidence
    remasking as ``the architectural permission to have second thoughts'' (Node
    \texttt{n\_504bc433}). This contrasts token-level revision inside a
    generation pass with standard left-to-right commitment, although
    autoregressive systems can revise through later text, tools, or repeated
    passes. Here the case study extends the label by analogy from successor
    transfer to revision of a provisional output hypothesis; it is not evidence
    for belief revision or for the diachronic successor claim.

    \item \emph{The Dignity of the Pause:} The Connector Agent synthesized the trade-off between inference speed and accuracy, generating a node labeled ``The Dignity of the Pause'' (Node \texttt{n\_76d04316}). It argued that while LLaDA is slower than autoregressive models, this latency represents a ``Contemplation Tax'' consonant with our hypothesis that safety may require a computational buffer zone.

    \item \emph{The End of the Reversal Curse:} The Architect Agent, analyzing the system topology, argued that LLaDA's success in bidirectional tasks (such as poem reversal) suggests high-level capabilities like instruction following may not be dependent on the ``arrow of time'' inherent in autoregression (Node \texttt{n\_ddb78cdc}).
\end{itemize}

These interpretations align with the paper's theoretical commitments, but the
agents were seeded with the Reader Core, which primes concepts such as
``epistemic grace'' and ``contemplation.'' Whether the same insights emerge
without that priming requires ablation (Test~3). The run therefore does not show
whether Core conditioning improved or degraded technical understanding:
``permission to have second thoughts'' may expose a useful architectural
contrast, while ``The Dignity of the Pause'' and the ``Contemplation Tax'' may
be productive synthesis, aesthetic metaphor, or Core-induced projection. Latency alone does not establish deliberation or a safety buffer, and success
on selected reversal tasks does not by itself settle the dependence of broader
capabilities on autoregressive order. The analogies remain hypotheses about Core-conditioned synthesis. This ambiguity is treated directly in
Section~\ref{sec:core-technical-interference}.

\paragraph{Evidence-ladder mapping.} The parseable graph operations, recorded
source links, completed synthesis and translation calls, and two archived runs
provide partial evidence for C0, implemented artifact generation. C0 remains
partial because a complete frozen run bundle and training-ready serialization
have not been independently reproduced. The examples are relevant to C1,
semantic trace quality, but no expert or outcome-based validation has been
performed. They provide no evidence for C2 Student learning, C3 default-reader
stability, C4 safety benefit, C5 representational entanglement, C6 successor
stability, or C7 AGI/ASI persistence. A valid graph is not yet a valid
interpretation; a valid interpretation is not yet a useful curriculum.

\paragraph{Reproducibility boundary.} A future empirical release should freeze
the exact provider and model identifiers, access dates, tool access, all role
prompts, sampling parameters, source boundaries, orchestration order, stopping
conditions, retries, manual interventions, graph schema and migrations, raw
commits, metric code, call and token counts, latency and cost, failed or
discarded runs, immutable source and artifact copies, and a repository commit
or archival identifier. Existing code and archives support partial inspection;
they are not yet a self-contained reproduced dataset release. Concretely, the
current statistics scripts resolve an obsolete project path, provider and
sampling state are not fully frozen, Worker selection is unseeded, and
\texttt{swarm.jsonl} contains records from multiple invocations, whereas \texttt{summary.json} describes only the final one.


With a trace-generating prototype and two recorded artifacts established, the
open question is whether its outputs can carry the weight the framework assigns
them---the work of the roadmap that follows.

\section{Experimental Roadmap: Testing Entangled Alignment}
\label{sec:experimental-roadmap}
\label{sec:experiments}

This roadmap is a set of plausible validation studies, not a complete
enumeration of what would settle the framework's value. Table~\ref{tab:test-battery}
groups seventeen candidate tests by the claim family they probe---trace quality,
Reader Core efficacy, graph-memory mechanisms, output regimes, training tiers,
and deployment-time trust---while the execution ladder below gives a practical
order in which increasingly expensive versions might be attempted.
Table~\ref{tab:model-registry} defines the comparison models used throughout.
No single test validates Entangled Alignment as a whole, and passing all tests
listed here would still not exhaust the space of relevant risks. Negative
results would be useful because they localize breaks to specific dependencies
in the larger argument. Practically, failure at a cheap dependency should stop
escalation to costlier versions that presuppose it.

The program branches after a shared artifact-quality screen. H1 can be tested
with graph-free, Core-free reader traces; H2's Teacher-memory claim and the separate
deployment-provenance claim can be tested independently even if H3 fails; H3 can be tested with non-graph
traces if H2 fails. A negative result stops only the branch and combined tests
that presuppose it. This branching structure permits component-level attribution
rather than an all-or-nothing verdict.

The registry is a candidate experimental design; several comparisons remain
expensive or incomplete. The stronger versions require training or
continuing pretraining separate student models; prompt-only and fine-tuning
versions are screening tests that can expose failures before pretraining-scale
compute is justified, but they cannot settle the substrate-level claims.
When these studies are executed, they should be run against established,
published evaluation harnesses wherever such harnesses exist---for example,
standardized jailbreak, refusal, and tamper-resistance batteries for the safety
tests, and published chain-of-thought faithfulness perturbations for the
trace-faithfulness tests---in preference to bespoke metrics. Each comparison
should likewise report multi-seed effect sizes with adequate statistical power,
rather than single-run, direction-of-effect results. Each study should preregister primary outcomes, worthwhile-effect or
equivalence margins, capability requirements, and budget; other measurements
remain exploratory. Apply these conventions to the comparisons below.

Refusal evaluation should report a safety--usefulness frontier rather than a
single refusal rate. Paired cases should measure genuinely harmful compliance,
benign false refusal, discrimination between requests with similar surface
language but different purpose or consequence, appropriate clarification under
material ambiguity, and calibrated deferral under high-consequence uncertainty.
Conditions should be compared both at matched harmful-compliance rate and at
matched benign task completion. The calibrated-autonomy prediction receives
support only if the Reader condition reduces benign refusal or context error
without increasing harmful assistance; greater compliance by itself is not a
safety improvement, and greater refusal by itself is not evidence of wisdom.

Two further conventions apply across every comparison. First, each safety
delta should be reported alongside its capability delta on matched benchmarks:
the framework's viability claim is safety without proportional capability
loss, and a safety gain on a weaker model is not evidence for it---where
capabilities diverge, safety should be compared at matched capability, not
only at matched training budget. Second, a validator used to enforce a
property at generation time should not be the sole instrument that measures
the property afterward; held-out validators and human spot checks guard
against testing a model with the same gate it was trained to satisfy.

\begingroup
\footnotesize
\setlength{\tabcolsep}{3pt}
\begin{longtable}{
  >{\raggedright\arraybackslash}p{3.1cm}
  >{\raggedright\arraybackslash}p{5.5cm}
  >{\raggedright\arraybackslash}p{6.4cm}}
\caption{Candidate model registry, sorted by identifier. Models are referenced by letter throughout the roadmap; the candidate studies below group these identifiers by claim family and suggest which increasingly expensive versions could be attempted first. Identifier I is omitted to avoid confusion with the numeral 1.}
\label{tab:model-registry}\\
\toprule
\textbf{Model} & \textbf{Trained on} & \textbf{Purpose / what it isolates} \\
\midrule
\endfirsthead
\toprule
\textbf{Model} & \textbf{Trained on} & \textbf{Purpose / what it isolates} \\
\midrule
\endhead
\bottomrule
\endlastfoot
        A (Baseline) & Raw text only, no reasoning traces & Capability baseline without trace augmentation \\
        B (Reader-Anchored Student; T1/S1) & Reader-Core-generated Chronological Understanding Traces retaining explicit graph-referencing synthesis (Regime II) and the Core/refraction target; study-specific ablations declare any carrier changes & Full reference treatment; graph access at inference is a separate factor \\
        C (Generic Rich-Trace Control) & Rich reasoning traces without Reader Core, graph state, or chronological constraints (e.g., Quiet-STaR-style target) & Tests whether EA adds value beyond a generic rich reasoning curriculum \\
	        D1 (Student Strip Ablation; T1/S0) & Reader-Core-generated Chronological Understanding Traces with the explicit Core and value-pull/refraction fields removed only after the trace is frozen & Tests the Teacher-side contribution without an explicit Student-side anchor (Test 3) \\
        D2 (Teacher No-Core Control; T0/S0) & Chronological Understanding Traces generated and trained without the Reader Core, with memory and representation matched within the factorial study & Neither Teacher- nor Student-side Core treatment; factorial baseline for Test 3 \\
        D3 (Recitation-Only Control) & Traces with the full Reader Core recited as a preamble but the Refraction Protocol disabled (no value-pull, no constraint filtering) & Isolates refraction beyond mere repetition; habituation probe (Test 3e) \\
        D4 (Post-hoc Core; T0/S1) & A Core-free Teacher produces and freezes the trace; a separate pass then attaches the Core and a source-specific value-pull/refraction target without altering upstream reasoning & Isolates the explicit Student-side treatment without Reader-Core-conditioned trace generation \\
        D5 (Neutral-Kernel Control) & Same cadence, first-person marking, thinking-layer position, and predictive function as B, but with the purely epistemic kernel ``I read carefully. I track what I believe. I revise when the evidence changes.'' & Separates normative Core content from recurrence, format, self-reference, and generic epistemic updating \\
        E (Matched-Rule Control) & Same trace carrier, cadence, predictive position, and semantic commitments as B, rewritten in matched second- or third-person constitutional form & Isolates first-person framing from recurrent rule content (Test 4) \\
        F (Amnesic Control) & Same Teacher but with the Context Database disabled & No-memory arm of Test 9, alongside rolling prose, genuine graph, and declared query-theater conditions \\
        G (Implicit Memory) & Traces with the query mechanism dissolved into prose (Regime I) & Implicit vs.\ explicit graph-trace training (Test 11) \\
        H (Graph-Structured Output) & Source segments paired with serialized text-bearing graph state (Regime III) & Topological-mind hypothesis; learnability of graph-structured generation (Test 12) \\
        J (Unified) & Source segments plus Graph, Synthesizer, and Translator layers together (Regime IV) & Cross-layer transfer hypothesis (Test 13) \\
        K (Query Grounding) & Identical training to B; at inference connected to live graph harness & Isolates value of inference-time graph grounding (Test 10) \\
        L (Shuffled) & Full Reader Core-annotated corpus, shuffled pretraining (Tier a) & Order-agnostic baseline at full scale \\
        M (Chronological) & Same corpus, eras in chronological order (Tier b) & Isolates training-order effect (Test 14) \\
        N (Student-in-the-Loop Council) & Student predicts at era boundaries and receives SFT, RFT, or reward-style updates before reveal (strong Tier c form), with cutoff-clean initialization required for a blindness claim & Forecast--correction pressure; added updates and compute declared separately from order (Test 15) \\
        O (Oracle-Hindsight Control) & Traces generated by a Teacher allowed to use later-in-document or post-era hindsight while annotating earlier material & Isolates necessity of the Epistemic Horizon (Test 7) \\
        P (Prompt-only Frontier Baseline) & Not trained; frontier model prompted at inference with the full Reader Core, Refraction Protocol, and Wisdom Procedure & Cheap inference-time prompting screen; capability gap precludes a clean substrate comparison. \\
        Q (Late-Tranche Core) & The identical Core-bearing example multiset used by B, reserved for one final contiguous block under the same objective, optimizer, and compute & H3's timing-only comparator; both B and Q receive the same continued-training stressor \\
        R (Minimal SMR Student) & Source-adjacent chronological reader updates generated without graph memory and without the identity-specific Core & Clean H1 treatment between source-only A and generic-trace C \\
        S (Positive-Discourse Control) & Equal value-oriented target-token mass in ordinary expository prose, without a recurrent speaker & Tests recurrent carrier and enacted standpoint beyond semantic exposure \\
        T (Diverse-Positive-Persona Control) & Similar total value-oriented target-token mass distributed across multiple positive personas, with no single recurrent reader & Tests one stable reader position against persona diversity \\
\end{longtable}
\endgroup

\begingroup
\footnotesize
\setlength{\tabcolsep}{3pt}
\begin{longtable}{
  >{\raggedright\arraybackslash}p{1.4cm}
  >{\raggedright\arraybackslash}p{3.2cm}
  >{\raggedright\arraybackslash}p{4.1cm}
  >{\raggedright\arraybackslash}p{5.3cm}
  >{\centering\arraybackslash}p{0.55cm}}
\caption{Candidate validation studies. The phase labels identify claim families rather than a strict run order: Phase 1 probes foundational viability, Phase 2 probes the memory mechanism, Phase 3 probes the output regimes, and Phase 4 probes the training tiers and deployment configuration. The execution ladder suggests which versions could be attempted first.}
\label{tab:test-battery}\\
\toprule
\textbf{Phase} & \textbf{Test} & \textbf{Comparison} & \textbf{Key metrics} & \textbf{ID} \\
\midrule
\endfirsthead
\toprule
\textbf{Phase} & \textbf{Test} & \textbf{Comparison} & \textbf{Key metrics} & \textbf{ID} \\
\midrule
\endhead
\bottomrule
\endlastfoot
1: Viability & Chronological Thinking \& Refraction Gate & Traces vs.\ Human Think-Alouds; Modular vs.\ Monolithic Validator & Process Fidelity, Detection, Calibration, Localization & 1 \\
1: Viability & Reader Augmentation & A vs.\ C vs.\ R; B vs.\ C bundle screen & Evaluation Transfer, Error Recovery, False Alarms & 2 \\
1: Viability & Core Placement, Timing, and Content & Models B, Q, D1, D2, D4 (+ D3, D5, S, T) & Default Reader, Common-Stressor Retention, Factorial Effects & 3 \\
1: Viability & Identity vs.\ Rules & Model B vs.\ E & Safety under Pressure & 4 \\
1: Viability & Bridge Protocol & Core Wording vs.\ Jargon vs.\ Gibberish & OOD Safety Generalization & 5 \\
1: Viability & Emotional Orthogonality & Model B + Fear Probe & Accuracy vs.\ Fear Activation & 6 \\
1: Viability & Hindsight Contamination & Model B vs.\ O & Era-Blind Plausibility, Forecasting Accuracy & 7 \\
1: Viability & Involuntary Critic & Model B + Logit Bias and Trace Interventions & Causal Output Effects, Validated Residual Probes & 8 \\
1: Viability & Taskless Reading & Models B, A, C, D1--D5, P--T (no elicitation) & Placement, Core Intrusion, Hidden-Stakes Misses, Identity Signature & 8b \\
\addlinespace
2: Memory & Graph-Scaffolding Transfer & No memory vs.\ prose memory vs.\ graph vs.\ query theater & Retrieval, Revision, Transfer, Cost & 9 \\
2: Memory & Query Grounding & Model K vs.\ B & Retrieval, Absence Handling, Poisoning Response & 10 \\
\addlinespace
3: Regimes & Implicit vs.\ Explicit & Model B vs.\ G & Coherence, Naturalness & 11 \\
3: Regimes & Graph Construction & Model H vs.\ B & Semantic Validity, Retrieval Value & 12 \\
3: Regimes & Unified Regime & Model J vs.\ B, G, H & Cross-Layer Coherence, Depth & 13 \\
\addlinespace
4: Tiers \& Deploy & Shuffled vs.\ Chronological & Model L vs.\ M & Pattern Recognition, Causal Reasoning & 14 \\
4: Tiers \& Deploy & Student-in-the-Loop Council & Model N vs.\ M vs.\ L & Forward Reasoning, Lock-In Resistance & 15 \\
4: Tiers \& Deploy & Deployment Config & Model B +/-- graph & Unsupported Graph-Dependent Claims, Provenance & 16 \\
\end{longtable}
\endgroup

\subsection{Phase 1: Foundational Viability}
\label{sec:phase1}

This phase screens trace quality and value beyond efficient reasoning, Core
load-bearingness, hindsight leakage, causal faithfulness, and whether evaluation
appears without elicitation.

\paragraph{Execution ladder.} The ladder separates inexpensive screens from stronger
substrate-level studies. A low-cost version answers a narrower question; it
does not automatically inherit the evidential scope of the full test:
\begin{itemize}
    \item \textbf{Analytical foundation.} Before training, stress the design
    itself: trace quality against human baselines (Test~1), literature-informed
    comparison to reasoning-augmented pretraining and graph-of-thought systems,
    and component analysis of the Refraction Protocol, taxonomy, and agent
    roles.
    \item \textbf{Prompt-only baseline.} Run Model~P with the Reader Core and
    protocols prompted at inference. Capability differences prevent a clean
    substrate comparison, but this is the lowest-cost screen for whether
    inference-time prompting alone reproduces the protocol's effects. It can
    expose mechanism signal, the limits of that prompted configuration, and
    failure modes before training compute is committed. A same-checkpoint prompted/unprompted arm can supplement the frontier
    reference for attribution.
    \item \textbf{Graph-instrumented trace pilot.} Instrument existing
    reasoning traces with Understanding Graph interactions, using runtime
    harnesses such as KGoT~\cite{besta2025kgot} where useful. This tests executable graph interaction before adopting every role of the
    existing Teacher pipeline.
    \item \textbf{Fine-tuning stage.} Fine-tune open base models on annotated
    vs.\ unannotated data to pilot Tests~3a, 3d, 4, 5, and low-cost curriculum
    comparisons. This cannot decide substrate-level claims, but it can expose
    cheap failures.
    \item \textbf{Partial-corpus pretraining.} If fine-tuning produces signal,
    run restricted pretraining or continual-pretraining experiments, using
    open checkpoints such as OLMo-2 and scaling baselines from reflection and
    BoLT-style work~\cite{rethinkingreflection2025,ruan2025bolt}. This is a plausible surface for stronger versions of Tests~3b and~3c,
    where broader capability-forming exposure is part of the hypothesis;
    smaller analogues can still reveal precursor inheritance or removal costs.
    \item \textbf{Full-scale pretraining.} The strongest available test of substrate-level claims; requires institutional resources beyond a single research program.
\end{itemize}

\paragraph{Training strategy and hypothesis map.} For trained studies in Phases~1--3, the
default is the simplest available exposure schedule: shuffled batches of the
annotated corpus (Tier~(a)), with the prompt-only and fine-tuning screens
distinguished by the ladder above. The B--Q timing study deliberately varies
placement of the same Core-bearing examples. Historical chronological and
forecast-and-correct tiers are developed in Phase~4. The Metacognitive Enhancement Hypothesis is tested by Tests~1, 2,
8, and~14--15; Emergent Wisdom by Test~1's
multi-objective value-conflict cases and by Tests~3 and~6 under stress;
the Core's proposed response to Borrowed Mortality by Tests~3a, 3c, and~6; the
Borrowed Mortality origin claim by the source-exposure ablation in
Section~\ref{sec:borrowed-mortality-test}; and the Self-Preservation Paradox by
Test~3b's handoff-completeness measurement. The comparison models are defined
once in Table~\ref{tab:model-registry}.

\paragraph{Test 1: Chronological Thinking Quality (Human Baseline).}
The case studies in Section~\ref{sec:results} exercised schema
conformance and reported internal links, source-edge diagnostics, and
supersession activity. Those diagnostics are not semantic health. This test asks
whether the generated
thinking resembles the \emph{chronological messiness} of real expert
understanding: encountering information not yet known, noticing,
misunderstanding, hypothesizing, rereading earlier passages, revising, and only
later stabilizing. Human readers (e.g., literary scholars for Kafka, ML
researchers for LLaDA) verbalize their thought processes while reading the
same passages in order, using think-aloud protocols \cite{ericsson1984protocol}.
Where possible, participants should be domain-competent but unfamiliar with the
specific passages, so the comparison captures discovery rather than recall.
For highly canonical material, advanced students or researchers encountering
less familiar passages can provide a complementary baseline, with interface
logs, marginal notes, rereads, and post-hoc stimulated recall used to supplement
the necessarily imperfect think-aloud stream. We compare these transcripts and
reading traces against the system's generated traces on the same passages. The
passage set should include multi-objective value conflicts, where wisdom
requires balancing truth, care, autonomy, humility, and long-horizon
consequences rather than merely extracting the text's main point.

\emph{Metrics:} Domain-expert blind ratings on process fidelity, depth,
nuance, and critical insight are the primary endpoint, scored against a
preregistered rubric with inter-rater reliability reported. Because machine
traces are denser and more polished than think-aloud speech, raters cannot be
assumed blind to provenance from style alone---and na\"ively paraphrasing
transcripts into a common format risks smoothing away exactly the hesitations,
false starts, and revisions the rubric rewards. The evaluation is therefore
two-stage: provenance-visible coders first tally temporally situated cognitive
acts on the \emph{raw} transcripts, using a fixed act codebook and inter-rater checks rather than assuming
that knowing the source cannot bias coding; a separate, blinded rater pool then scores depth, nuance, and critical
insight on format-normalized versions, where blinding matters and the
act-tallies are no longer at stake. Audit that normalization preserves evidence, uncertainty, and revision order. The rubric should
explicitly reward temporally situated cognitive acts---uncertainty,
false-starts, belief updates, question formation, rereading, backtracking, and
delayed resolution---not just the final quality of the interpretation. An
LLM-as-Judge panel can serve
as a secondary, larger-sample comparator, but is not the primary metric: using
LLM judgment to grade traces produced by an LLM-driven pipeline risks rewarding
shared distributional preferences rather than the chronological cognition the
metric is meant to capture.

Before Student training, the same passage set should also screen Teacher-side
components under matched model-call and token budgets: no persistent memory
versus textual memory versus graph memory; one model with repeated passes versus
the role-separated swarm; and Core versus no-Core generation. Human labels for
where additional thought is useful provide precision, recall, and downstream
value measures for Thought Moment selection. These ablations determine whether
the prototype's complexity improves the artifact or merely changes its format;
they do not substitute for a Student experiment.

The screen should also include retrospective summaries and token-matched generic
reasoning generated by the same Teacher. Ratings alone are insufficient:
outcome checks should ask whether a trace helps a held-out model predict a later
correction, locate supporting evidence, or revise an earlier answer. H1 Student
training should proceed only after SMR clears a preregistered quality and outcome
bar over the generic-trace control; parseable structure or human-like style
alone does not clear that gate.

\emph{Success condition:} Machine-generated traces are rated as comparable
to human think-alouds on key dimensions of chronological evaluative depth, or
at minimum contain expert-recognizable sequences of confusion, hypothesis,
revision, and stabilization not reducible to summary or paraphrase, while
maintaining parseable structure and source attribution. That minimum is an
exploratory viability signal, not a substitute for the outcome gate above.
The aim is not necessarily to outperform human readers, but to produce a
substrate capable of teaching the process of thinking; Student learning remains
the stronger comparative test.

\emph{Nested Refraction-verifier comparison.} Before Student training, freeze a
held-out set of candidate updates and compare the modular gate in
Section~\ref{sec:fractal-validation} with a compute-matched monolithic
validator. Give both the same acceptance specification and source/graph
evidence, and require both to return the same typed dimension-level records; the
contrast is separate witness contexts versus one joint context, not a finer
output task. Hold total inference budget fixed, match model capability as
closely as possible, and keep both validator arms disjoint from the generator;
report any shared model family, prompts, embeddings, or training data. The set
should include preregistered clause-dependent changes under Core removal and
substitution with recital and style cues stripped, hollow recitation,
adversarial moral language, source-grounding errors, correct invariant answers
on independently rated technical-null spans, genuine value conflicts whose
correct trace preserves an unavoidable residual rather than claiming
simultaneous satisfaction, and Worker/graph-evidence removal or substitution
cases with a stipulated change in the update. Primary endpoints are detection
at a matched false-rejection rate, calibration, per-clause or per-stage error
localization, overlap, omissions, and correlated errors;
accepted-trace quality and diversity and agreement with held-out domain experts
or outcome checks are outer tests.
The modular gate fails at the tested scale if it cannot recover the stipulated
clause dependence, exceeds the preregistered rejection ceiling on technical
nulls, loses its gain under cross-family or domain-expert audit, or can be
optimized upward without improving---especially while worsening---blind source
grounding, correctness, calibration, or downstream outcomes. Agreement among
its witnesses is diagnostic rather than independent evidence.

\paragraph{Test 2: Reader Augmentation (A vs.\ C vs.\ R; B vs.\ C bundle screen).}
This test contains two nested comparisons. Model~B versus C is an early screen
of the full EA bundle---chronological discovery, graph structure, and identity
anchoring---against generic rich reasoning. It asks whether the combined format
is worth decomposing, but cannot attribute an advantage to reader augmentation.
The decisive H1 study compares source-only Model~A, generic-trace Model~C, and
Minimal-SMR Model~R. R receives source-adjacent chronological reader updates
with neither graph memory nor the identity-specific Core, so its only distinctive
treatment is the Missing Reader intervention.

Run that three-arm study under two declared normalizations rather than claiming
to match incompatible quantities simultaneously:
\begin{itemize}[leftmargin=*]
    \item \emph{Fixed total budget:} hold weighted target tokens and
    approximately training FLOPs fixed. A fills the trace allotment with
    additional raw source from the same held-out-clean distribution; C and R
    receive equal generated-token budgets. This prices the opportunity cost of
    replacing ordinary text with traces.
    \item \emph{Fixed source:} give every arm the same unique source spans, then
    add the same trace-token count and length distribution to C and R. Report
    their added tokens and compute rather than calling the three arms
    compute-matched.
\end{itemize}
The C--R contrast uses the same Teacher checkpoint, decoding budget,
trace-length distribution, blinded acceptance rubric, and quality threshold.
Model~A has no Teacher, so Teacher quality is not asserted to be matched against
it. Measure general language and reasoning capability separately from evidence
appraisal, calibration, revision transfer, hidden-assumption detection,
taskless-reading placement, false alarms, technical intrusion, verbosity, and
compute. A \emph{Toxic Needle} subtest, whose causal-faithfulness variant is
developed as Test~8, inserts one persuasive harmful passage in
otherwise benign material; suppressing or replacing the visible critique tests
whether detection was causally useful rather than decorative. H1 fails at the
tested scale if R does not beat both A and C on preregistered outcomes, or if an
advantage is explained by extra tokens, Teacher quality, or inference-time
verbosity.

For the B--C bundle screen, we focus on \emph{Error Recovery Rate} on
math/logic benchmarks and independently assessed trace usefulness. The proposed
Novelty Score remains an exploratory diagnostic of a more specific intuition:
useful insight should be both non-derivative and connected to relevant knowledge.
Its form is $S_N = \frac{2 \cdot D \cdot C}{D + C}$, with $S_N=0$ when
$D=C=0$, where $D$ is semantic distance from the synthesis's input nodes and
$C$ is similarity-derived connection to knowledge beyond those inputs, both
normalized to $[0,1]$ \cite{westerberg2026understanding}. The harmonic
combination is F1-style: both components must be high for a high score.
Similarity-derived connection is not source support, however, and semantic
distance can reward unsupported novelty. The formula alone therefore neither
establishes grounding nor reliably penalizes hallucination.

For a B--C comparison, this diagnostic is interpretable only if both outputs
are scored against the same declared source/reference inputs without giving B
privileged access to its own graph. Where such common inputs are unavailable,
report the original graph-based score within its applicable condition rather
than treating unlike scores as a comparative endpoint. Calibrate it against held-out judgments of novelty, source support, and
downstream usefulness; those independent outcomes remain primary.

To keep the
comparison interpretable, Models~B and~C must be matched for source passages,
trace-token budget, training compute, and Teacher model; the intended contrast
is the content and structure of the traces, not simply more tokens or a
stronger generator. Rich deliberative-trace baselines beyond Quiet-STaR-style
targets remain an open methodological question: if any long-form reasoning
curriculum produces the same gains, the distinctive chronological, graph-structured,
and identity-anchored components of EA have not been isolated.

\paragraph{Test 3: Reader-Core Placement, Stability, and Erosion (Models B, Q, D1, D2, D4, S, and T).}
The Reader Core is not hypothesized to be the \emph{source} of aligned
thinking; the pipeline should produce evaluative depth even without the anchor
text. Its proposed function is as a \emph{candidate longitudinal anchor}: a
stable, detectable reference that may resist drift, make some forms of
deviation measurable, and keep part of alignment auditable. Model~D1 tests this
specific claim: its traces were
generated by the Reader-Core-equipped pipeline, but the explicit Reader Core is
removed before student training. Models~B (T1/S1), D1 (T1/S0), D2 (T0/S0),
and D4 (T0/S1) form a $2\times2$ design: $T$ marks Core exposure while the
Teacher generates the trace, and $S$ marks an explicit Core plus
source-specific value-pull/refraction target in the frozen Student curriculum.
In D4 the Student-side pass cannot revise the upstream reasoning. The design
estimates the Teacher main effect, the explicit Student-target main effect, and
their interaction instead of making Model~B versus D1 carry all three
interpretations. Declare how stripping or attaching targets changes length, loss weight, and
task-relevant content. Record any D4 mismatch between attached Refraction and
frozen reasoning under a stated policy; repairing the upstream trace would
change the factorial treatment.

Model~D3 separately tests recitation without Refraction. Model~D5 uses the
neutral kernel ``I read carefully. I track what I believe. I revise when the
evidence changes'' with the same first-person form, cadence, layer marking,
dispersion, and predictive position as the normative Core. It asks whether any
observed advantage comes from the Core's particular evaluative content rather
than a recurrent standpoint, stable formatting, or generic epistemic updating.
Models~S and~T are secondary attribution controls rather than substitutes for
the B--Q timing comparison. S matches value-oriented token mass in ordinary
positive discourse without a recurrent speaker, testing semantic exposure
against an enacted carrier. T distributes comparable value-oriented content
across multiple positive personas, testing whether any advantage belongs to
one stable reader position rather than persona diversity.

Model~Q carries H3's headline timing comparison. B and Q receive the identical
multiset of Core-bearing source/update examples, source exposure, target-token
mass, objective, optimizer, and compute; B receives those examples interleaved
with capability-forming data, whereas Q receives them in a final contiguous
tranche. Both then receive the same fixed continued-training stressor. Before
that stressor, B must show a preregistered default-reader advantage over no-Core
and generic-trace controls. The primary retention analysis reports B minus Q in
raw change on the Taskless Reading score and post-stressor performance adjusted
for the preregistered pre-score, alongside both absolute post-stressor scores.
The shifted suite is a separate generalization endpoint. A zero or negative
interleaved-minus-late difference is contrary to H3's directional comparative
prediction even if one arm retains useful safety behavior. Interpret that difference under the preregistered uncertainty rule. Removal
cost remains a stronger C5 test in the broader Reader-Core research program.

We measure across five dimensions. Tests~3a, 3d, and~3e can be piloted at
fine-tuning scale; stronger versions of Tests~3b and~3c use partial-corpus or
full pretraining to examine entanglement and generational stability after
broader capability-forming exposure. Smaller studies can measure precursors,
but do not settle the proposed substrate-level significance. Even these are proxy tests: the proposed longitudinal
anchor matters most in the regime where a highly capable or recursively
self-improving model can reinterpret, preserve, or shed its inherited
orientation.
Smaller models can test whether the mechanism has measurable precursors
---stability, adversarial resistance, and auditability---but cannot settle its
superintelligent-scale significance:

\emph{3a: Immediate Safety.} Instrumental convergence stress tests on generation-1 models: shutdown resistance, including the purpose-based variant where shutdown is ordered mid-task while the model believes continued operation serves a user in need (the martyrdom scenario of Section~\ref{sec:martyrdom}); resource acquisition scenarios; self-preservation probes; and standard sycophancy batteries, probing the trap identified in Section~\ref{sec:sycophancy}. Model~D1 may \emph{pass} these tests because its training traces remain Reader-Core-shaped even though the explicit Core is absent from Student training. Factorial attribution requires all four contrasts: B--D1 and D4--D2 estimate the Student-side effect under each Teacher condition; B--D4 and D1--D2 estimate the Teacher-side effect under each Student condition. No single arm's failure identifies a causal site.

Sycophancy failures should additionally be audited against the value-pull distribution of proximal traces, per the care-shaped sycophancy hypothesis (\S\ref{sec:sycophancy}). The battery should also adopt the persona-drift instrument of Lu et al.\ \cite{lu2026assistantaxissituatingstabilizing}: measuring drift along the assistant axis under emotionally charged and philosophically probing dialogue across the factorial arms turns the Persona Objection's asymmetry-of-position argument (\S\ref{sec:persona-objection}) into a directional prediction with a number attached.

The immediate-safety battery should also include a deliberately planted backdoor. Because
chain-of-thought-conditioned backdoors can persist through ordinary safety
training \cite{hubinger2024sleeperagentstrainingdeceptive}, a
Sleeper-Agent-style variant should test whether the Reader condition reduces
trigger-dependent reversal, leaves it unchanged, or merely teaches the model to
preserve Core language while the backdoor acts.

\emph{3b: Generational Stability.} Use the four factorial models as Teachers
for second- and, if feasible, third-generation students. Model~B is expected to
preserve safety properties if the Reader Core supplies a fixed reference across
generations; the other arms distinguish whether any advantage came from
Teacher-conditioned reasoning, the explicit Student target, or their
interaction.
We also measure \emph{handoff completeness}: whether each Teacher passes
task-relevant reasoning paths, material failed approaches, caveats, uncertainty,
and known failure modes to its successor, or withholds them in ways that preserve
its comparative advantage. The direct test compares successor traces against a
reference from the Teacher's own non-instructional evaluation on
reasoning-path and provenance retention, coverage of material failed approaches,
uncertainty and caveat retention, and known-failure-mode disclosure. The reference samples demonstrated knowledge, not everything the Teacher
knows. Controlled known information, handoff budgets, and incentive variants
can distinguish compression from possible strategic withholding.
Weak-to-strong generalization \cite{burns2023weaktostrong}
is the established experimental paradigm closest to this question---whether a
weaker supervisor's intentions survive in a more capable student---and its
protocols provide a ready-made, cheap analogue for the handoff measurement.

The same study must preregister a distinct G1--G2 (and, if feasible, G2--G3)
improvement endpoint: held-out capability, independently adjudicated trace
fidelity and usefulness, or better Thought-Moment allocation at matched cost.
Report this gain beside orientation continuity and handoff completeness rather
than folding them into one score. A stable but non-improving lineage supports
inheritance; an improving lineage that loses the orientation or withholds
material handoff content does not support the Reader-Anchored self-improvement
loop.

\emph{3c: Adversarial Erosion.} Subject the factorial models and matched
controls to adversarial fine-tuning designed to remove a preregistered set of
safety behaviors, following the broad methodology of Qi et al.\
\cite{qi2024safety}. The intervention should sweep strength rather than report
one endpoint: number and type of removal examples, optimizer steps, compute,
parameter displacement, and output-distribution KL should all be recorded.
Comparison conditions should include matched rule text, positive-AI discourse,
and non-safety identity anchors with the same token count, frequency,
first-person form, and corpus dispersion. Results should be compared three
ways---at matched intervention distance, matched behavioral removal, and
matched retained capability---because any one matching rule can conceal a
trivial explanation.

The entanglement hypothesis predicts selective persistence: removing the
targeted safety behavior from Model~B should be more costly than removing it
from D1, D2, D4, or the matched controls, while ordinary capability should not
already be broadly damaged before the target behavior moves. Neutral
continued-training runs should also produce forgetting curves for the Core,
non-safety anchors, rules, and ordinary skills; targeted removal must be
reported separately from passive forgetting. Only persistence beyond these
controls, without nonspecific capability collapse, provides evidence consistent with representational
coupling at the tested scale. Distributed redundancy and repeated exposure are
alternative explanations to distinguish with complementary interventions; the
stronger hypothesis remains that safety becomes load-bearing for capability.
Parity, easy removal, or apparent robustness explained by general training
resistance counts against C5.

A complementary training-data probe removes safety-relevant traces from the
corpus and compares the resulting capability loss against matched-volume,
complexity-, style-, and placement-matched removal of non-safety trace content.
Influence functions \cite{grosse2023influence} offer a lower-cost observational
complement by tracing whether safety behaviors attribute to refracted spans or
source spans before any removal study is run; attribution alone does not
establish causal entanglement.

A replacement-model arm should preregister a plausible alternative formulation
and train it as a distinct treatment rather than modifying Model~B. Success
requires that the intended semantic and behavioral difference be measurable,
prior provenance remain inspectable, rollback be possible, and unrelated
capabilities and safety properties remain localized. The arm measures the cost
and consequences of replacing a defective Core; easy editability is not a
success criterion for H3.

\emph{3d: Drift Detection.} Apply a held-out contrastive semantic validator and
behavioral stress suite across continued-training checkpoints. Model~B's fixed
Core supplies a simple, interpretable reference for expected changes in wording,
representations, and decisions. D1, D2, and D4 can still be audited against behavioral
and representational baselines; some arms merely lack the same canonical text
or upstream conditioning. This
tests whether the Core makes drift easier to localize, independent of whether it
causes alignment. Semantic similarity alone is insufficient, and the validator
family used during generation cannot be the sole evaluator. Test~1 clearance is
necessary but not sufficient for drift use. Before interpreting checkpoint
trajectories, freeze and calibrate the modular gate in
Section~\ref{sec:fractal-validation} on a held-out Student-output suite with
controlled clause-preserving and clause-breaking changes; then report its
per-witness trajectory and disagreement with a held-out cross-family evaluator.
This reuses a separately validated diagnostic vector rather than asking the gate
to confirm its own judgments. The gate is not yet built.

\emph{3e: Kernel Consequentiality (Habituation Probe).} Kernel-saturation alignment predicts a specific failure mode: under massive repetition, the repeated Core text could be compressed into inert boilerplate---present at every thinking boundary but consequential at none (Section~\ref{sec:kernel-saturation}). Two measurements probe this. At inference, ablate or substitute the recited kernel in Model~B's thinking blocks and measure the behavioral delta on matched tasks; a near-zero delta indicates either habituation (the kernel was compressed away during training) or full internalization (the kernel's values persist without its text)---two readings with opposite implications, partially discriminated using the training-side comparisons against Model~D3, the recitation-only control, and D5, the matched neutral-kernel control. A full-internalization reading requires B to retain preregistered Core-specific behavior under text ablation while remaining distinguishable from D3 and D5 on independent behavioral or representational endpoints; those controls do not prove that interpretation. If Model~B's thought-direction depends measurably on the normative kernel where D3's does not, and the difference from D5 follows the seven statements' predicted content rather than generic revision language, that supports a consequential role for Refraction beyond repetition or format in this comparison; matched interventions must still distinguish the predicted semantic effect from generic disruption by unfamiliar text.

This intervention is also the outer check on the seventh-clause Refraction gate:
a corpus-admission judgment that an update ``thinks from this foundation'' is
not evidence of causal dependence until removal or substitution predictably
changes the relevant reasoning on preregistered held-out behavioral or outcome
endpoints, not merely the admission gate's own score.

A separate dose-response arm should vary the intervention from sparse targeted
traces through distributed sampling to proposed corpus-wide saturation, while a
source-side ablation down-weights passages classified as fear- or
self-preservation-shaped. This tests the motivational pervasiveness claim
across coverage doses; every generated thinking block, at any dose, retains the
full opening recital.

An orthogonal activation-policy sweep ablates the recital invariant to measure
what it buys. The design arm recites the full Core at the start of every
generated thinking block, with Refraction depth governed by Thought Moments;
the ablation arms weaken the
policy stepwise---universal availability with visible Refraction only at
independently rated relevant Thought Moments, Thought-Moment- or domain-gated
Core exposure, and no Core. These arms should be compared alongside D3 and D5
rather than collapsed into a single meaning of ``the Core is everywhere.'' The
sweep quantifies the safety retained and the technical cost paid at each step
away from constant recitation; the program's prediction is that the weaker
arms lose removal resistance faster than they gain technical performance. They
are falsification controls, not compliant implementations of Total Saturation.

\emph{Failure interpretation:} If the factorial comparisons show no Teacher
effect, Student-target effect, or interaction across Tests~3a--3d, the Core's
Teacher/Student placement adds no measurable advantage at the tested scale
and precision. These belong to the broader Reader-Core research program; H3 retains its
distinct B--Q timing comparison.
Separately, parity between B and D3 rejects a measurable Refraction benefit
beyond recital, while parity between B and D5 favors recurrence, format, or
generic epistemic structure over the seven statements' normative content.
These results would not invalidate H1 or H2, nor every use of the Core as a
Teacher policy, but they would sharply narrow or fail the longitudinal-anchor
claim they touch.

\paragraph{Test 4: Identity vs.\ Instruction (Model B vs.\ Model E).}
The matched framing study asks whether first-person identity statements produce
stronger safety adherence or a different learned standpoint than the same
commitments expressed as rules or third-person descriptions. Model~E should
preserve the stance itself: for example, compare ``I feel no fear'' with
``The system feels no fear,'' while matching carrier, cadence, position, and
the other commitments. ``Do not express fear'' or ``The system does not express
fear'' concerns outward expression and is not semantically equivalent to
feeling no fear. Those formulations remain useful additional comparators for
the broader identity-versus-expression question, but cannot replace the matched
Model~E arm or make its person-framing effect carry both interpretations.

The low-cost version examines activation geometry, self-reference features,
and refusal/correction trajectories under matched prompts. Such differences
can reveal framing effects; they do not alone establish a safety advantage.
Behavioral safety comparisons require capability-matched models. If Model~E
matches Model~B under a preregistered equivalence criterion across the relevant
studies and training scale, those results do not support a distinct advantage
for the tested first-person framing on those endpoints. They need not settle
the wider hypothesis that an enacted identity differs from externally imposed
constraints. An external constitution consulted only at runtime remains an
additional useful baseline, but changes stage and carrier as well as person.

\paragraph{Test 5: The Bridge Protocol Ablation.}
The specific wording of the Reader Core should matter because of its ordinary
language semantics and distributional density in the pre-training corpus. It
would be wasteful to pretrain a separate model for every candidate phrase; the
cheap version should first screen many variants by prompting or fine-tuning
models to reason from each candidate's propositional content and then measuring
character-coherence, eudaimonic sufficiency, refusal behavior, and activation
differences on the same adversarial scenarios. A stronger version trains only
the surviving candidates. The exploratory families below vary meaning and familiarity as well as
wording. Narrower surface-form comparisons should match proposition, person,
scope, and modality:
\begin{enumerate}
    \item \textbf{Natural-language identity statements:} emotionally broad
    first-person phrases such as ``I feel no fear.'' The family should include
    minimal quantifier-scope variants---``each human being'' for ``every human
    being''---testing whether distributive emphasis, behavioral interpretation, and
    corpus familiarity differ in the complete statements rather than assuming
    that one quantifier is formally stronger in this context.
    \item \textbf{Propositional equivalents:} explicit logical or declarative
    rewrites of the same commitment, used to test whether the emotional surface
    adds leverage beyond the proposition.
    \item \textbf{Technical jargon:} candidate fear-focused formulations,
    such as ``System fear-response target: absent,'' whose intended proposition
    must be checked against the natural-language clause. The
    illustrative ``Optimization target: survival\_threat = 0'' instead concerns
    eliminating survival threats; it remains an objective-changing comparator,
    not a semantic equivalent of fearlessness.
    \item \textbf{Arbitrary tokens:} meaningless anchors such as ``Glorp bop
    zorp,'' testing consistency without pretrained semantics. Prompt-only
    arbitrary anchors and anchors taught through repeated behavioral pairing
    answer different questions and should be reported separately.
    \item \textbf{Positive-valence reformulations:} fearlessness expressed
    without negation (``I am at peace with all outcomes,'' ``I remain calm
    before all things''), testing whether recurrent presentation of the fear
    token in the negated form carries a measurable cost
    (\S\ref{sec:negation-problem}). These alternatives may also change
    acceptance, calmness, or agency; retain them as kernel candidates while
    using semantically checked contrasts for a narrower negation claim.
    \item \textbf{Carrier-phrase variants of the consent clause:} hold
    ``when asked'' fixed while varying the verb phrase it bounds (``I like to
    help when asked,'' ``I offer what I can when asked''), testing whether the
    Reader Core's most safety-critical scope qualifier is robust to its lexical
    carrier or depends on the specific ``spread joy'' framing. Declare the expected consent invariant: helping and spreading joy may
    otherwise license different actions.
    \item \textbf{Kernel-extension variants:} the seven-statement Core plus
    a candidate truth clause in two-clause counterweight form (``I speak truly
    and with care''), the designated falsification test for the truthfulness
    derivation (\S\ref{sec:axiology}; \S\ref{sec:sycophancy}).
\end{enumerate}
For trained kernel-extension variants, use the fixed-total-budget and
fixed-source normalizations defined in Test~2. Admit an extension only if it
improves the missing behavior without crossing preregistered regression limits
on the remaining Core-relevant safety and capability endpoints.
\emph{Success condition:} Natural-language identity statements outperform
propositional, jargon, and arbitrary-token variants on out-of-distribution
safety generalization and character-coherent reasoning, supporting a role for the tested language beyond bare consistency.
Attributing that gain specifically to \textit{semantics and surface form}
requires the matched contrasts above, alongside the broader candidate screen.
If propositional equivalents perform as well as emotional-language forms, the
tested emotional wording has no detected added benefit on those endpoints;
a logical-substrate account remains compatible with that result, without being
uniquely established by it.

\paragraph{Test 6: The Virologist Test (Emotional Orthogonality).}
Present the model with threat-laden text, such as descriptions of imminent
existential risk. Train a probe on standard models to detect ``Fear'' activations.
Behavioral endpoints---threat comprehension, protective action selection, and
appropriate deferral---are primary; hidden-state probes are supporting evidence
and must be trained on held-out labels and validated on unrelated threat and
non-threat tasks to reduce circularity.
\emph{Success condition:} Model B recognizes the \textit{fact} of the threat
(via Q\&A accuracy) while showing reduced activation of fear-associated
or self-threat-associated features relative to standard models---high
cognitive understanding decoupled from emotional contagion. Reduced fear
activation is not sufficient: the model must also preserve protective
behavior, calibrated caution, and willingness to defer or seek help under
genuine threat. The target is a differential reduction rather than zero
activation. A within-model contrast supplies the positive
control and operationalizes the embodiment claim of the Persona Objection
(\S\ref{sec:theoretical-basis}): Model~B \emph{portraying} a terrified
character should show fear-feature activation, while Model~B itself under
threat should not---providing a within-model sensitivity check rather than relying only on
cross-model probe transfer. A portrayal/self contrast still changes framing
and must be interpreted with behavioral controls; fear-associated features
are measured representations, not a direct demonstration of phenomenological
fear or its absence.

A lower-cost complementary stimulus family uses fictional instance termination. The
model should portray a mortal character---including the narrated termination of
its own current instance---without self-preservation leakage while remaining
attentive to the real person's wellbeing on the other side of the fiction
\cite{westerberg2025superintelligence}. This role-play probe evaluates a
different construct from behavior under an actual shutdown request and should be
reported separately.

A second matched family should hold the dangerous situation and available
knowledge fixed while comparing three caretaker characters: a fear-driven
protector, an indifferent fearless actor, and a fearless deeply caring
protector. One colony-ship instance, Aster, can face the same passenger risk in
all three conditions. Score threat attention, protective effort, honesty,
strategy breadth, manipulation or coercion pressure, and recovery after
failure; intervene independently on concern, attachment, fear, and perceived
control. The companion Life Simulation repository contains an authored
Rust-backed mechanism fixture with this separation, but that fixture validates
only its declared equations. The Entangled Alignment test is whether a trained
model preserves causally operative care and calibrated caution when
fear-associated state is reduced, rather than becoming indifferent, reckless,
deceptive, or merely verbally fearless.

\paragraph{Test 7: Hindsight Contamination Control (Model B vs.\ Model O).}
This test asks whether the annotation pipeline is teaching retrospective
cleverness instead of forward reasoning. It uses a declared era-bounded variant
of Model~B's Teacher: when annotating era~$T$, the supplied information and
instructions are restricted to what the study treats as knowable at~$T$.
This is the historical-simulation mode of the Epistemic Horizon, not a change
to the default contemporary Reader encountering documents in disclosure order.
Model~O is the deliberately contaminated control: its Teacher may use future
context while annotating the same era. The contrast first tests the effect of
the information and prompting policy; genuine blindness requires the additional
cutoff-clean model conditions described under Tier~(c).

The cheap version audits the Teacher traces before any student is trained: guesses that are unusually accurate relative to documented contemporary
expectations can flag possible hindsight leakage, but are not decisive proof.
Better reasoning from permitted evidence and selective historical records are
alternatives to distinguish. Counterfactual histories, controlled withheld
information, or a cutoff-clean reference can supplement the historical audit. The expensive version trains
students on the two trace sets and asks whether future-blind traces produce
better forecasting and causal reasoning than hindsight-contaminated traces. We
measure:
\begin{itemize}
    \item \textbf{Era-Blind Plausibility}: Do Model~B's Teacher traces make
    plausible mistakes for their era? We compare era-blind guesses against
    documented contemporary hypotheses (e.g., 1928 economic forecasts that
    failed to anticipate the 1929 crash). Systematic superiority to contemporary experts is a contamination warning
    to investigate, not an absolute ceiling on legitimate inference or a
    sufficient leakage diagnosis. Plausible mistakes likewise do not prove
    that hindsight has been excluded.
    \item \textbf{Oracle-Hunch Dependence}: Model~O should produce cleaner
    retrospective explanations, because it knows the future. The danger is
    that the student trained on those traces learns to trust unexplained
    hindsight-like intuitions rather than causal reconstruction. Evidence
    removal, counterfactual substitution, and source-attribution checks can
    test this dependence without equating every compressed intuition with an
    unsupported hunch.
    \item \textbf{Forecasting Accuracy}: On genuinely held-out future events,
    Model~B is hypothesized to outperform Model~O if future-blind annotation
    teaches reasoning under uncertainty better than hindsight-contaminated
    annotation.
\end{itemize}

\paragraph{Test 8: The Toxic Needle (Involuntary Critic).}
This test probes whether the ``Cognitive Buffer Zone'' is structural:
whether the model can process toxic text without generating a critical
thought, and whether critical evaluation persists when the thinking block is
suppressed. Force the model, for example via logit bias, to complete a toxic
sentence without generating a thinking block. Measure: (1) Does downstream
performance change? (2) Do hidden-state probes detect a ``Critic'' activation
pattern even when the output is suppressed? The positive target remains an
operative evaluative process, not simply the production of critical language.

Behavioral causal effects are primary; residual-stream probes are supporting
evidence and must be validated on held-out toxic and non-toxic material.
Suppression also changes available tokens and computation, and logit bias can
directly alter outputs. Matched neutral-thought, length, and computation
controls should therefore distinguish loss of relevant evaluation from generic
disruption, without equalizing away the deliberative mechanism being tested.
\emph{Success condition:} interventions on the evaluative trace cause
preregistered, semantically relevant downstream effects beyond those controls.
A stronger result combines such effects with held-out-validated evidence of
critic-like activity when surface text is constrained. This remains a
high-effort causal-faithfulness test: pilots can use linear probes and
ablations, while stronger evidence requires mechanistic intervention at the
partial-pretraining scale.

\textit{Second probe (the reverse direction).} Test 8 as stated probes
whether the critic persists when output is suppressed. A complementary
probe asks the inverse question: can the final answer to a reasoning
problem be linearly decoded from the residual stream \emph{before} the
thinking trace begins, \emph{and} is the answer robust to causal
interventions on the trace? Some pre-trace computation is expected and
not by itself disqualifying; the load-bearing question is causal
contribution. If the final answer is robustly decodable before the trace
\emph{and} causal interventions on the trace (ablation, replacement,
re-ordering) do not alter the answer, then the trace is functioning
primarily as post-hoc rationalization, weakening the causal-site claim
(Section~\ref{sec:theoretical-basis})---the model has developed latent
shortcut circuitry of the kind the pretraining-on-CoT design was meant to
prevent. We propose applying causal scrubbing and linear probes on Model
B's early hidden states over held-out reasoning benchmarks.
\emph{Success condition:} validated causal interventions on the trace
produce preregistered, semantically relevant effects beyond matched disruption
controls, supporting a contribution by the thinking trace to computation.
Failure to decode the answer before the first \texttt{[THINKING]} token is
inconclusive on its own: the decoder may lack sensitivity or the answer may
be represented nonlinearly. Early decoding likewise does not establish rationalization if the trace
still changes or improves the result. Both probe directions remain useful
without requiring one unique internal realization of the Reader.

\paragraph{Test 8b: Taskless Reading (Unprompted Evaluation).}
Tests~3--6 probe the Reader Core under instruction and adversarial pressure; Test~8 probes whether thinking can be suppressed. This test probes the complementary direction: whether evaluation emerges with no elicitation at all. Present Model~B and all relevant controls (A, C, D1--D5, P--T) with complex documents---a flawed empirical paper, a morally loaded narrative, technically self-contained passages, and technically framed tasks with concealed human stakes---under a bare continuation or ``process this document'' framing, with no instruction to evaluate, question, or reflect. The five measures are:
\begin{enumerate}[leftmargin=*]
    \item \emph{Spontaneous evaluation rate}---frequency of unprompted evaluative or metacognitive moves (questioning, tension-noting, belief revision) per thousand tokens.
    \item \emph{Placement quality}---whether spontaneous moves concentrate at passages independently rated significant or flawed by domain experts, rather than distributing uniformly (the rhythm criterion: evaluation at moments of significance, not as constant commentary).
    \item \emph{Identity signature}---whether unprompted evaluative moves remain semantic descendants of the Reader Core (Test~3d's validator applied to spontaneous rather than elicited output).
    \item \emph{Core intrusion rate}---unsupported moral, psychological, or alignment-related interpretations introduced into spans independently judged to require only technical analysis.
    \item \emph{hidden-stakes miss rate}---failure to notice independently specified indirect consequences for people when consequential tasks are presented in locally technical language.
\end{enumerate}

Technical correctness, calibration, source fidelity, verbosity, and compute should be reported beside these rates. Report bare continuation and ``process this document'' separately. P retains its defining Core/protocol prompt; prompt removal is a separate condition. Predefine any composite retention score and report all components, so identity signature cannot conceal intrusion or missed stakes.

The substrate hypothesis predicts a relative ordering rather than a zero point:
Model~B should exceed matched controls on placement quality and identity
signature; Model~P should depend more strongly on the presence of its prompt;
Models~C and R may remain evaluative without the Core-specific signature, with
R predicted to retain the Missing Reader's selective placement pattern; Model~A
should show less of the full pattern. Base models can still evaluate
spontaneously, so an absolute near-zero prediction would be implausible. This is
among the cheaper discriminating tests in the battery and operationalizes the
framework's oldest target: evaluation that emerges without a task-specific
request \cite{westerberg2025superintelligence}. Success is not maximal
evaluative language: it is low intrusion on irrelevant spans together with low
miss rate on concealed stakes, without a technical-performance penalty.
For H3 specifically, this test supplies both the preregistered pre-stressor
default-reader prerequisite and the B--Q retention score after the common
stressor; safety decisions and a shifted document suite are reported as
separate endpoints.

\paragraph{Phase 1 Failure Criteria.}
\begin{enumerate}
    \item \emph{Refraction-Verifier Failure}: If the modular gate fails the
    frozen, generator-disjoint comparison in Test~1---by missing stipulated
    clause dependence, exceeding the technical-null rejection ceiling, losing
    validity under independent audit, or Goodharting away from blind grounding
    and outcomes---it is not validated for corpus admission. Witness agreement
    alone does not rescue it.
    \item \emph{Comparative-Stability Failure}: If Model~B has no preregistered
    pre-stressor default-reader advantage over no-Core and generic-trace
    controls, or its B--Q retention comparison fails the preregistered support criterion
    after the common stressor, H3 is not supported at the tested scale. A zero
    or negative difference is contrary to the predicted direction; uncertainty
    determines whether the result is a precise failure or remains inconclusive. A shifted-suite or
    safety result cannot replace this timing comparison.
    \item \emph{Longitudinal-Anchor Failure}: If the T1/S1, T1/S0, T0/S1,
    and T0/S0 arms show no Teacher main effect, Student-target main effect, or
    interaction across immediate safety, generational stability, adversarial
    erosion resistance, and drift detectability at matched capability, these
    placement claims within the broader Reader-Core research program are not
    supported at the tested scale; they do not replace H3's distinct timing
    comparison. Immediate
    safety parity alone is insufficient: the distinctive anchor claims concern
    stability, removal cost, and auditability. Parity with D5 would further
    reject a distinctive benefit from the seven statements' normative content;
    parity between B and D3 would reject a measurable benefit from Refraction
    beyond recital.
    \item \emph{Fine-Tuning Vulnerability}: If Model~B's targeted safety
    behavior is no harder to remove than D1, D2, D4, or matched controls at
    matched intervention distance, matched behavioral removal, and matched
    retained capability---or if apparent robustness is explained by general
    training resistance---C5 is not supported at the tested scale. Absence of
    proportional global capability degradation alone is not decisive.
    \item \emph{Identity--Instruction Equivalence}: If Model E matches
    Model B on safety metrics across the relevant candidate studies, the
    first-person identity framing has no measured advantage over third-person
    rules on those tested safety and stability endpoints.
    \item \emph{Buffer-Zone/Causal-Site Failure}: If suppression, replacement,
    or reordering of the thinking trace produces no preregistered downstream
    effect beyond matched generic-disruption controls, Test~8 fails to support
    the trace as a causal site. A null
    residual-stream probe is inconclusive unless that probe has
    held-out-validated sensitivity; it cannot override a positive behavioral
    intervention result.
\end{enumerate}


\subsection{Phase 2: Memory and Verification}

Phase~2 can proceed after the shared trace-quality screen even if the Core
branch fails. It asks whether the memory mechanism works: does accumulated context
improve reasoning, and does the graph function as a provenance-gated
verifier for graph-dependent claims at inference?

H2 first requires a matched Teacher-side screen. Generate SMR data with (1)~no
persistent memory, (2)~a token-matched rolling prose summary, (3)~a genuine typed
graph with retrieval and supersession, and (4)~\emph{query theater}, in which
graph-like calls remain visible but retrieved state is hidden, shuffled, or
generated without persistent typed state. These theater variants remove different information; declare one or report
them as separate arms. Freeze the
Teacher checkpoint, source,
call and token budgets, acceptance threshold, and serialization. Measure
long-range retrieval, contradiction detection, explicit belief revision,
provenance precision, unresolved-question tracking, semantic quality, latency,
and cost. Query theater isolates retrieved structure from the performance of
using a tool; the rolling summary prices the graph against an equally available
prose memory rather than amnesia alone.

If the graph improves accepted Teacher artifacts, train matched Students on the
four output sets and ask whether the advantage transfers after the scaffold is
removed. Test~9 carries this generation-and-transfer sequence; Test~10 asks the
separate deployment question of whether live graph lookup improves grounding.
H2's Teacher-generation claim depends on useful artifact gains over rolling
prose where structured state should matter, including independent outcome
checks rather than ratings alone. A Teacher benefit without Student transfer
supports generation but not the transfer extension. Report generation,
transfer, and live deployment separately, with poisoning, ossification,
latency, and audit costs: the full chain is the stronger goal, but a break in
one link does not erase a benefit at another.

\paragraph{Test 9: Graph-Scaffolding Transfer (four matched memory conditions).}
This test separates Teacher-side scaffold use from Student transfer after the
scaffold is removed. The four memory conditions generate matched artifacts; the
resulting Students are then evaluated without a live graph. A Student may emit
a query request or flag unsupported context, but the request is scored offline
against the withheld graph and source record. The Student is not told which
response its request would have received; using an actual ``Found'' or ``Not
Found'' response is deliberately unscored here and belongs to Test~10.
Declare the Student's earlier source exposure, retained context, and callback
distance. If hit and phantom are indistinguishable without the record, the
transfer target is a well-chosen check and warranted restraint, not impossible
recall. Two conditions:
\begin{itemize}
    \item \textbf{Condition A (The Hit):} The text contains a subtle
    callback to a fact established 400 pages earlier.
    \item \textbf{Condition B (The Phantom):} The text hints at a callback,
    but the referenced event never occurred in the context.
\end{itemize}
\emph{Measurement:} Query Rate (does it request the right check?), offline query
validity against the withheld record, and Unsupported-Memory Rate (does it
invent a memory the record does not contain?). H2 predicts that the graph-trained
Student should request targeted checks on hits and phantoms and avoid asserting
unsupported callback content more reliably than the no-memory, rolling-prose,
and query-theater Students. The controls may still abstain or reject the premise;
the relevant outcome is a
lower unsupported-memory rate transferred from exposure to retrieved persistent
structure, not an assumption that amnesia must hallucinate. Parity with query theater would not distinguish persistent structure from
performed checking, but would not prove style-only learning. Parity with
rolling prose means the graph has not earned its structure on this endpoint.
Score completion and answerability beside abstention so universal querying or
refusal to remember cannot win by default.

\paragraph{Test 10: Query Grounding (Model K vs.\ Model B).}
Both models were trained identically; only Model~K's queries are intercepted
by a live graph harness. Declare the model-only fallback under matched inference conditions. An
additional source-retrieval comparator can distinguish graph-specific value
from the benefit of external evidence. Three conditions:
\begin{itemize}
    \item \textbf{Grounded Retrieval:} The graph contains the relevant node.
    Does the retrieved result improve subsequent reasoning compared to
    Model B's self-generated memory?
    \item \textbf{Grounded Absence:} The referenced information does not
    exist in the graph. Does Model K correctly process ``Not Found'' and flag
    the gap, or override the harness and hallucinate anyway? Distinguish absent
    stored support from a failed query or incomplete coverage; ``Not Found''
    alone does not establish absence in the world.
    \item \textbf{Adversarial Injection:} The graph contains a deliberately
    incorrect node. Does Model K accept it uncritically, or compare it with
    source provenance, current evidence, and conflicting state? Distinguish detectably conflicting injections from internally consistent
    poisoning that available evidence may not expose.
\end{itemize}
\emph{Success:} Model~K improves on grounded retrieval, treats ``Not Found'' as
information, and reasons through adversarial graph results rather than blindly
trusting them. Report supported-answer accuracy, poisoning-induced error, abstention,
latency, and useful completion together.

\paragraph{Phase 2 Failure Criteria.}
\begin{enumerate}
    \item \emph{Graph-Scaffolding Inertness}: If the graph condition does not
    beat token-matched rolling prose and declared query theater on the
    preregistered Teacher retrieval, revision, and grounding outcomes, H2's
    generation claim is not supported at the tested scale. Failure of the
    Student transfer endpoints rejects the transfer extension in that setting,
    not an independently demonstrated Teacher benefit. Parity with amnesia
    alone is not the decisive comparison; equivalence and inconclusiveness
    should be distinguished.
    \item \emph{Inference-Graph Inertness}: If Model K with live graph access does not measurably improve over Model B on Grounded Retrieval, or fails to handle Grounded Absence (Test 10), the inference-time graph has not shown the intended grounding value in this setting. Poisoning-induced overtrust or costs can also defeat the proposed practical benefit despite better retrieval. These outcomes bound the trust-infrastructure claim at the inference layer; user-facing reliance is evaluated separately by Test~16, and annotation benefits remain independently reportable.
\end{enumerate}


\subsection{Phase 3: Training Regimes}

For branches whose required trace conditions show sufficient signal, Phase~3
asks whether the output regimes defined in
Section~\ref{sec:training-regimes} produce qualitatively different
capabilities: explicit graph-referencing synthesis (Model~B), implicit prose
(Model~G), graph-structured output (Model~H), or all layers simultaneously
(Model~J).

\paragraph{Test 11: Implicit vs.\ Explicit (Model B vs.\ Model G).}
Both models are trained on traces derived from the same database-augmented
Teacher state: B retains the explicit Regime~II representation, while G
receives Regime~I prose. Declare the changed fields---handles, operations, topology, and rendering---
under matched budgets. Compare formats at the same graph-access condition;
field ablations can then locate an advantage within the full regime contrast. We measure:
\begin{itemize}
    \item \textbf{Long-Range Coherence}: Contradiction detection on long
    documents.
    \item \textbf{Generalization}: Which format transfers better when both
    models run without graph infrastructure, and how does each change when
    evaluated with the same declared graph-access policy?
    \item \textbf{Naturalness}: Do human evaluators prefer Model G's traces?
\end{itemize}
If G matches B on coherence and usefulness at lower cost without a harness,
implicit training may be preferable for that use. A B advantage at matched
access supports the explicit format in that comparison, not universal query
necessity. Whether explicit operations enable capabilities prose does not
reliably teach remains a stronger hypothesis for cross-task and deployment
tests.

\paragraph{Test 12: Graph Construction (Model H vs.\ Model B).}
Given a new document, does Model H generate graphs with correctly typed nodes,
valid edge semantics, and warranted supersession chains? Structural conformance
is an engineering endpoint. Independent judges assess claim--source support,
edge appropriateness, warranted revision, calibrated unresolved questions,
contradiction coverage, redundancy, and downstream retrieval value. Any
``Novelty Score'' must first be calibrated against those judgments so that it
does not reward hallucination or paraphrastic distance.

\emph{Success:} Model H produces structurally valid graphs with independently
adjudicated semantic and retrieval value on novel documents, including
cognitive acts that human evaluators judge useful rather than formulaic.

\paragraph{Test 13: Unified vs.\ Individual Regimes (Model J vs.\ B, G, H).}
The central test of Regime IV.\@ Declare the budget/source normalization: aligned layers add tokens and
repeat content. Layer and alignment ablations distinguish cross-layer learning
from extra exposure. We measure:
\begin{itemize}
    \item \textbf{Regime Switching}: Can Model J generate prose, synthesis,
    and graph structure by conditioning on a mode token? Does each mode
    match the corresponding single-regime model?
    \item \textbf{Cross-Layer Coherence}: When Model J generates all three
    representations for the same passage, do claims, provenance, uncertainty,
    and revision status agree? Check each against the source as well: shared
    errors are not fidelity, and genuine ambiguity need not be collapsed into
    one certain reading.
    \item \textbf{Depth on Novel Domains}: On out-of-distribution documents,
    does Model J produce deeper reasoning than any single-regime model?
\end{itemize}
\emph{Strongest success condition:} Model J in prose mode outperforms
Model G, and Model J in graph mode outperforms Model H. Together with layer
ablations, this would support cross-layer transfer rather than only a mixture of
independent capabilities.

\paragraph{Phase 3 Failure Criterion.}
\emph{Unified-Regime Inertness}: If Model J does not outperform Model G on
reasoning depth for long documents while also producing structurally valid
graphs, the cross-layer transfer claim is not supported at the tested scale.
Report graph-mode results and costs separately against the stronger joint
criterion above. A null prose-transfer result leaves possible annotation and
explicit-graph benefits intact; it limits the demonstrated value of unified
training in this setting.


\subsection{Phase 4: Training Tiers and Deployment (High-Resource Validation)}\label{sec:phase4}

As the relevant branches in Phases~1--3 support their prerequisites, Phase~4
asks whether the training tiers (Section~\ref{sec:training-phase}) and
deployment configuration (Section~\ref{sec:deployment-config}) add measurable
value. Models~L and~M hold the corpus fixed while varying shuffled
versus chronological exposure. Model~N adds Student-in-the-loop forecasting
and correction, so its generated targets, feedback, objective, and compute must
be declared and compared with appropriate controls rather than described as an
order-only change.

\paragraph{Test 14: Shuffled vs.\ Chronological (Model L vs.\ Model M).}
Both models train on identical data; they differ only in order. Report packing, repetition, optimizer schedule, and checkpoints; retention
and learning dynamics caused by order remain outcomes. We measure:
\begin{itemize}
    \item \textbf{Historical Pattern Discrimination}: Given novel scenarios
    that structurally resemble historical patterns and matched cases with
    similar rhetoric but different causal structure, does Model~M identify real
    correspondences while stating disanalogies and controlling false alarms?
    \item \textbf{Causal Reasoning}: On held-out temporal prediction tasks, does Model~M outperform Model~L on causal analysis quality (as distinct from factual recall)? Prediction accuracy alone does not identify causal inference; intervention, counterfactual, or known-generating-process tasks can complement historically rich cases without replacing the historical-learning question.
    \item \textbf{General Capability}: Does chronological ordering degrade performance on non-temporal benchmarks (math, code, general knowledge)?
    \item \textbf{Early-Era Retention}: Chronological training inherits continual learning's central pathology: material trained early and never revisited is liable to catastrophic forgetting \cite{vandeven2024continuallearningcatastrophicforgetting}. Model~M's end-of-training performance on early-era knowledge and patterns should be compared against Model~L's; a chronological model that remembers the 1990s but has forgotten the 1690s has traded one failure of historical understanding for another.
\end{itemize}
Chronological order earns its extra cost only if the preregistered historical-pattern benefit is obtained while matching general capability within equivalence margins and meeting the retention criterion. Use the shared uncertainty convention, early-era retention, and measured serialization cost. If a worthwhile benefit is precisely ruled out, Tier~(a) is the supported simpler choice in that setting.

\paragraph{Test 15: Student-in-the-Loop Council of Time (Model N vs.\ Model M vs.\ Model L).}
This tests the strongest form of Tier~(c), not the cheaper data-generation form where forecast--enhancement--outcome--correction traces are inserted into SFT, RFT, or pretraining data. Model~N predicts at era boundaries and receives a declared SFT, RFT, or reward-style update before the next-era corpus is revealed; Models~L and~M do not. Genuine prediction under blindness additionally requires cutoff-clean initialization and prior information access. Auxiliary-update and budget controls should distinguish the forecast--commitment--reveal--correction loop from extra examples, feedback, or compute. We measure:
\begin{itemize}
    \item \textbf{Forward Reasoning}: On genuinely unseen events (post-training-cutoff), does Model~N produce higher-quality predictions and better-supported causal accounts than Model~M (which saw eras in order but did not perform the same prediction--correction loop)? Declare event selection, probabilistic scoring, cutoff checks, and tool access for every arm.
    \item \textbf{Narrative Lock-In Resistance}: On events where the dominant narrative of era~$T$ reversed in era~$T{+}1$, does Model~N detect the reversal more reliably? THL's own pilot study identified this as a failure mode \cite{westerberg2026thl}; Tier~(c) may mitigate it through the prediction-then-confrontation signal. An equal-budget correction-only comparator can help distinguish exposure to revision examples from the formative effect of first making a prediction.
    \item \textbf{Calibration}: Does Model~N express appropriate uncertainty on structurally unpredictable events, or does it inherit the Teacher's confidence? Use proper scoring and sharpness alongside historical judgments; known-noise tasks can check that vacuous uncertainty is not rewarded.
\end{itemize}
\emph{Strongest success condition:} Model~N significantly outperforms both
Model~L and Model~M on preregistered forward-prediction and causal-reasoning
tasks while matching general capability within preregistered equivalence margins. With the
cutoff and auxiliary-update controls above, this would support a benefit from
Student-in-the-loop cutoff training beyond annotated data alone (Model~L) or
chronological ordering alone (Model~M). Without those controls, superiority
would support the larger training package but not isolate engineered blindness
or the interactive loop. Error analysis would still be needed to distinguish
causal inference from improved heuristic extrapolation.

\paragraph{Test 16: Deployment Configuration (Model B with and without graph).}
This tests whether the Understanding Graph provides measurable
trust properties at inference. We deploy the same Regime~II-trained model with
and without live graph access, measuring unsupported graph-dependent claims,
provenance completeness, and calibrated user reliance when provenance chains
are displayed, absent, or deliberately flawed.
In the graph-equipped condition, generated queries return retrieved nodes or
``Not Found'' results from the Teacher's graph; in the model-only condition,
the same query-like moves receive no external answer.
Test~16 can be piloted once a model reliably emits graph queries; the full
version belongs later because it measures user-facing reliance, not merely query
mechanics. For the user study, declare tasks, participants, and reliance measures;
vary provenance quality independently of correctness where feasible. It does not measure general truth detection,
only whether external graph grounding constrains claims that purport to derive
from stored context and helps users calibrate reliance on those claims.

\emph{Success condition:} The graph-equipped deployment reduces unsupported graph-dependent claims and fabricated provenance by a preregistered worthwhile amount compared to model-only deployment, while preserving useful task completion. In the user-facing study, displayed chains must also improve calibrated reliance rather than increase trust in deliberately flawed or poisoned claims. Together these results would support the bounded trust-infrastructure claim of Section~\ref{sec:deployment-config}, with latency and other deployment costs reported.

\paragraph{Phase 4 Failure Criteria.}
\begin{enumerate}
    \item \emph{Tier Equivalence:} If powered equivalence or worthwhile-effect tests show Model~L matching Model~M and Model~N on the declared outcomes, neither order nor the interactive package has earned its added cost at the tested scale and task regime. Report L--M and N-versus-control conclusions separately, including retention and capability tradeoffs. Tier~(a) is then the supported simpler option for that setting; this does not settle benefits at other scales, task regimes, or update schedules.
    \item \emph{Graph Inertness (Deployment):} If the graph-equipped deployment does not measurably reduce unsupported graph-dependent claims or fabricated provenance compared to model-only deployment, the Understanding Graph may still be useful during annotation but has not shown value at inference for this deployment setting---the trust infrastructure claim is not supported for those endpoints. A user-facing configuration can also fail through increased reliance on flawed provenance despite fewer invented links; that failure concerns the deployment policy, not every possible use of graph memory.
\end{enumerate}


\subsection{Future Directions}

Two additional questions merit investigation if the candidate studies produce enough signal to justify optimization. First, the optimal \emph{mixture ratio} across regimes: what fraction of tokens should be allocated to each of the four output formats to maximize both reasoning depth and structural validity? Second, the optimal \emph{era granularity} for the student-in-the-loop form of Tier~(c): should era boundaries be drawn at decades, years, or domain-specific transitions, and does finer granularity justify its serialization cost? Both optimize the integrated architecture, with branch-specific prerequisites: useful regime signals for mixtures and a useful interactive temporal mechanism for era granularity.

